\section{Lattice Yang--Mills theory: the scaling limit and the mass gap}\label{ym:sec:lattice}
The problem of Jaffe and Witten \cite{JaffeWitten2006}: for a compact simple
gauge group $G$, construct a nontrivial quantum Yang--Mills theory on
$\mathbb R^4$ and show that its Hamiltonian $H$ has
\begin{equation}\label{ym:gap-statement}
 \operatorname{spec}(H)\subset\{0\}\cup[\Delta,\infty),\qquad\Delta>0 .
\end{equation}
Here $\operatorname{spec}(H)$ denotes the \emph{spectrum} of the self-adjoint
Hamiltonian $H$, including any continuous spectrum, with the vacuum energy
normalized to zero. The condition excludes spectral values in $(0,\Delta)$:
every nonzero excitation energy is at least $\Delta$.
A formal proof \cite{OpenAIYM2026}, checked by the Lean kernel and reviewed in
\cite{YMGroup2026}, constructs such a theory as a scaling limit of Wilson's
lattice gauge theory, in the sense of Theorem~\ref{ym:thm:existence}. It does
not assert Eq.~\eqref{ym:gap-statement}, and it does not assert that the limit
is unique. This section fixes the objects, derives the constants, states what
is proved, and reduces Eq.~\eqref{ym:gap-statement} to statements about
partition functions on finite tori. The statement of
Theorem~\ref{ym:thm:existence} follows the definition-by-definition account in
\cite{Armstrong2026}.

\subsection{Gauge group, inner product, Casimir constant}
\begin{definition}\label{ym:def:gauge-group}
A gauge group is a triple $(m,G,\mathfrak g)$: an integer $m\ge1$, a closed
subgroup $G\subset U(m)$, and the real subspace
\begin{equation}\label{ym:lie-algebra}
 \mathfrak g=\{X\in M_m(\mathbb C):\ \exp(sX)\in G\ \text{for all }s\in\mathbb R\}.
\end{equation}
$\mathfrak g$ is simple if it is nonabelian and its only ideals are $\{0\}$
and $\mathfrak g$.
\end{definition}
Every $X\in\mathfrak g$ is skew-Hermitian, $X^\dagger=-X$. On $\mathfrak g$ put
\begin{equation}\label{ym:inner}
 \langle X,Y\rangle_{\mathfrak g}=-\frac1m\operatorname{Re}\operatorname{tr}(XY),
 \qquad \|X\|_{\mathfrak g}^2=\langle X,X\rangle_{\mathfrak g}
 =\frac1m\operatorname{tr}(X^\dagger X)\ \ge0 .
\end{equation}
It is invariant: $\langle gXg^{-1},gYg^{-1}\rangle_{\mathfrak g}
=\langle X,Y\rangle_{\mathfrak g}$ for $g\in G$, by cyclicity of the trace.
With $\operatorname{ad}(X)Y=[X,Y]$ and the Killing form
$K(X,Y)=\operatorname{tr}_{\mathfrak g}(\operatorname{ad}(X)\operatorname{ad}(Y))$,
the Casimir constant $C_{\mathfrak g}$ of a simple $\mathfrak g$ is defined by
\begin{equation}\label{ym:casimir}
 K(X,Y)=-C_{\mathfrak g}\,\langle X,Y\rangle_{\mathfrak g}.
\end{equation}
Both sides are invariant symmetric bilinear forms on a simple Lie algebra, so
they are proportional; $K$ is negative definite on a compact simple
$\mathfrak g$, so $C_{\mathfrak g}>0$. If the inner product is rescaled,
\begin{equation}\label{ym:casimir-rescale}
 \langle\cdot,\cdot\rangle'=c\,\langle\cdot,\cdot\rangle_{\mathfrak g}
 \quad\Longrightarrow\quad C'=\frac{C_{\mathfrak g}}{c},
\end{equation}
because $K$ does not depend on the inner product.

\begin{proposition}\label{ym:prop:casimir-sun}
For $G=SU(N)\subset U(N)$, $N\ge2$: $K(X,Y)=2N\operatorname{tr}(XY)$ and
$C_{\mathfrak g}=2N^2$.
\end{proposition}
\begin{proof}
For $A,B\in M_N(\mathbb C)$ the map $Z\mapsto AZB$ on $M_N(\mathbb C)$ has
trace, in the basis of matrix units $E_{ij}$,
\begin{equation}\label{ym:trace-AZB}
 \sum_{i,j}(AE_{ij}B)_{ij}=\sum_{i,j}A_{ii}B_{jj}
 =\operatorname{tr}(A)\operatorname{tr}(B).
\end{equation}
Expand the double commutator,
\begin{equation}
 \operatorname{ad}(X)\operatorname{ad}(Y)Z=[X,[Y,Z]]=XYZ-XZY-YZX+ZYX,
\end{equation}
and apply Eq.~\eqref{ym:trace-AZB} to each term, with
$(A,B)=(XY,1),(X,Y),(Y,X),(1,YX)$ and $\operatorname{tr}(1)=N$:
\begin{equation}\label{ym:killing-gl}
 \operatorname{tr}_{M_N(\mathbb C)}\big(\operatorname{ad}(X)\operatorname{ad}(Y)\big)
 =N\operatorname{tr}(XY)-\operatorname{tr}(X)\operatorname{tr}(Y)
 -\operatorname{tr}(Y)\operatorname{tr}(X)+N\operatorname{tr}(YX)
 =2N\operatorname{tr}(XY)-2\operatorname{tr}(X)\operatorname{tr}(Y).
\end{equation}
$M_N(\mathbb C)=\mathfrak{sl}_N(\mathbb C)\oplus\mathbb C1$;
$\operatorname{ad}(X)\operatorname{ad}(Y)$ annihilates $1$ and preserves
$\mathfrak{sl}_N(\mathbb C)$, so the trace over $\mathfrak{sl}_N(\mathbb C)$
equals Eq.~\eqref{ym:killing-gl}. $\mathfrak{sl}_N(\mathbb C)$ is the
complexification of $\mathfrak{su}(N)$, so for $X,Y\in\mathfrak{su}(N)$ the
real trace over $\mathfrak{su}(N)$ equals the complex trace over
$\mathfrak{sl}_N(\mathbb C)$. With $\operatorname{tr}X=\operatorname{tr}Y=0$,
$K(X,Y)=2N\operatorname{tr}(XY)$. For skew-Hermitian $X,Y$,
$\overline{\operatorname{tr}(XY)}=\operatorname{tr}(Y^\dagger X^\dagger)
=\operatorname{tr}(YX)=\operatorname{tr}(XY)$, so $\operatorname{tr}(XY)$ is
real and Eqs.~\eqref{ym:inner}, \eqref{ym:casimir} give
\begin{equation}
 2N\operatorname{tr}(XY)=-C_{\mathfrak g}\Big(-\frac1N\operatorname{tr}(XY)\Big)
 =\frac{C_{\mathfrak g}}{N}\operatorname{tr}(XY),\qquad C_{\mathfrak g}=2N^2 .
\end{equation}
\end{proof}
So $C_{\mathfrak{su}(2)}=8$ and $C_{\mathfrak{su}(3)}=18$.

\subsection{Links, plaquettes, Wilson measure}\label{ym:sec:wilson-measure}
The finite-lattice formulation separates the space of configurations, the
action, and the measure used to integrate observables. Background for this
construction is given in \cite{GattringerLang2010}, Sections~2.2--2.4 and~3.1.

\subsubsection{Sites, spacing and oriented links}
Let $\ell\ge1$ and $\Lambda_\ell=(\mathbb Z/\ell\mathbb Z)^4$. A site $y$ is
a four-tuple of integer coordinates, each taken modulo $\ell$. With lattice
spacing $a$, its physical coordinate is $x=ay$, understood modulo the
physical period $L=a\ell$ in each direction. Thus $\ell$ counts sites in one direction;
$a$ carries length. The coordinate vectors $e_1,\dots,e_4$ specify lattice
directions and are distinct from the Lie-algebra basis used later.

A positive link $(y,\mu)$ connects $y$ to $y+e_\mu$. A configuration assigns
one matrix in $G\subset U(m)$ to each positive link:
\begin{equation}\label{ym:configuration-space}
 U\colon\Lambda_\ell\times\{1,2,3,4\}\longrightarrow G,
 \qquad \mathcal U_\ell=G^{4\ell^4}.
\end{equation}
Here $m$ is the matrix representation dimension, whereas $4\ell^4$ counts
the independently assigned positive links. The reverse link is determined
by the same matrix, with its origin shifted:
\begin{equation}\label{ym:reverse-link}
 U(y,-\mu):=U(y-e_\mu,\mu)^{-1}.
\end{equation}
In particular the reverse of the link starting at $y$ in direction $\mu$
starts at $y+e_\mu$, and its matrix is
$U(y+e_\mu,-\mu)=U(y,\mu)^{-1}$. There is no second independent matrix for
this reverse orientation.

Our transport convention is fixed by test vectors $\psi(y)\in\mathbb C^m$:
$U(y,\mu)\psi(y+e_\mu)$ is compared with $\psi(y)$. These test vectors do
not add a dynamical matter field to the configuration or action.
If $\psi^\gamma(y)=\gamma(y)\psi(y)$, covariance requires
\begin{align}
 U^\gamma(y,\mu)&=\gamma(y)U(y,\mu)\gamma(y+e_\mu)^{-1},
 \label{ym:link-gauge}\\
 U^\gamma(y,\mu)\psi^\gamma(y+e_\mu)
 &=\gamma(y)U(y,\mu)\gamma(y+e_\mu)^{-1}
   \gamma(y+e_\mu)\psi(y+e_\mu)\notag\\
 &=\gamma(y)U(y,\mu)\psi(y+e_\mu).
 \label{ym:link-transport}
\end{align}
This convention determines the order of multiplication around a path.

\subsubsection{Plaquettes and gauge transformations}
For $\mu<\nu$, traverse the four sides of an elementary square in the order
\begin{equation}\label{ym:plaquette-path}
 y\ \longrightarrow\ y+e_\mu\ \longrightarrow\ y+e_\mu+e_\nu
 \ \longrightarrow\ y+e_\nu\ \longrightarrow\ y.
\end{equation}
The matrices for the third and fourth steps follow from
Eq.~\eqref{ym:reverse-link}. The plaquette variable and Wilson energy are
\begin{align}
 U_{y;\mu\nu}&=U(y,\mu)\,U(y+e_\mu,\nu)\,U(y+e_\nu,\mu)^{-1}\,U(y,\nu)^{-1},
 \label{ym:plaquette}\\
 E_\ell(U)&=\sum_{y\in\Lambda_\ell}\ \sum_{1\le\mu<\nu\le4}
 \Big(1-\frac1m\operatorname{Re}\operatorname{tr}U_{y;\mu\nu}\Big).
 \label{ym:wilson-energy}
\end{align}
There are $\binom42=6$ choices of an unordered direction pair, hence
$6\ell^4$ plaquettes. For a unitary matrix $V$, its eigenvalues have unit
modulus, so $|\operatorname{tr}V|\le m$. Each energy term lies in $[0,2]$,
giving $0\le E_\ell(U)\le12\ell^4$.

Apply Eq.~\eqref{ym:link-gauge} to the four factors in
Eq.~\eqref{ym:plaquette}. The intermediate gauge matrices cancel:
\begin{align}
 U^\gamma_{y;\mu\nu}
 &=\gamma(y)U(y,\mu)\gamma(y+e_\mu)^{-1}\cdot
   \gamma(y+e_\mu)U(y+e_\mu,\nu)\gamma(y+e_\mu+e_\nu)^{-1}\notag\\
 &\quad\cdot\gamma(y+e_\nu+e_\mu)U(y+e_\nu,\mu)^{-1}\gamma(y+e_\nu)^{-1}\cdot
   \gamma(y+e_\nu)U(y,\nu)^{-1}\gamma(y)^{-1}\notag\\
 &=\gamma(y)\,U_{y;\mu\nu}\,\gamma(y)^{-1}.
 \label{ym:plaquette-gauge}
\end{align}
Thus the plaquette matrix is gauge covariant. Cyclicity of the trace gives
\begin{equation}\label{ym:plaquette-trace-invariant}
 \operatorname{tr}U^\gamma_{y;\mu\nu}
 =\operatorname{tr}\big(\gamma(y)U_{y;\mu\nu}\gamma(y)^{-1}\big)
 =\operatorname{tr}\big(U_{y;\mu\nu}\gamma(y)^{-1}\gamma(y)\big)
 =\operatorname{tr}U_{y;\mu\nu},
 \qquad E_\ell(U^\gamma)=E_\ell(U).
\end{equation}
The trace is invariant, although the matrix need not be unchanged.

\subsubsection{Haar measure, normalization and observables}
Let $h$ be normalized Haar measure on the compact group $G$. It has
$h(G)=1$ and is invariant under left and right multiplication. For an
integrable $f$ and fixed $g_L,g_R\in G$,
\begin{equation}\label{ym:haar-invariance}
 \int_G f(g_LVg_R^{-1})\,h(dV)=\int_G f(V)\,h(dV).
\end{equation}
Assign this measure independently to the positive links before inserting
the action, and write
\begin{equation}\label{ym:product-haar}
 \mathcal DU:=\prod_{(y,\mu)}h\big(dU(y,\mu)\big),
 \qquad \int_{\mathcal U_\ell}\mathcal DU=1.
\end{equation}

\begin{definition}\label{ym:def:wilson-measure}
For $\beta\in\mathbb R$, the Wilson measure is
\begin{align}
 \mu_{\ell,\beta}(dU)&=\frac{1}{Z_{\ell,\beta}}
             e^{-\beta E_\ell(U)}\,\mathcal DU,
 \label{ym:wilson-measure}\\
 Z_{\ell,\beta}&=\int_{\mathcal U_\ell}e^{-\beta E_\ell(U)}\,\mathcal DU.
 \label{ym:wilson-partition}
\end{align}
\end{definition}
For $\beta\ge0$, the energy bound gives
\begin{equation}\label{ym:finite-partition-bounds}
 0<e^{-12\beta\ell^4}\le e^{-\beta E_\ell(U)}\le1,
 \qquad e^{-12\beta\ell^4}\le Z_{\ell,\beta}\le1.
\end{equation}
Consequently $Z$ is a finite positive normalization. At $\beta=0$,
$Z_{\ell,0}=1$ and the links have independent Haar laws. For nonzero
$\beta$, the plaquette factors generally couple links; the product measure
$\mathcal DU$ is the integration measure, not a claim of independence in
the interacting law.

For a bounded observable $\mathcal O\colon\mathcal U_\ell\to\mathbb C$,
\begin{equation}\label{ym:observable-expectation}
 \langle\mathcal O\rangle_{\ell,\beta}
 =\frac{\int\mathcal O(U)e^{-\beta E_\ell(U)}\,\mathcal DU}
        {\int e^{-\beta E_\ell(U)}\,\mathcal DU},
 \qquad\langle1\rangle_{\ell,\beta}=1.
\end{equation}
For a fixed gauge transformation, define $\Gamma_\gamma(U)=U^\gamma$.
Each link is multiplied by fixed group elements on its two ends. Equation~\eqref{ym:haar-invariance} gives
$\mathcal D(U^\gamma)=\mathcal DU$, and
Eq.~\eqref{ym:plaquette-trace-invariant} gives
$E_\ell(U^\gamma)=E_\ell(U)$. Changing variables therefore yields
\begin{equation}\label{ym:measure-gauge-invariance}
 \int\mathcal O(U^\gamma)e^{-\beta E_\ell(U^\gamma)}\,\mathcal DU
 =\int\mathcal O(V)e^{-\beta E_\ell(V)}\,\mathcal DV,
 \qquad \langle\mathcal O\circ\Gamma_\gamma\rangle_{\ell,\beta}
 =\langle\mathcal O\rangle_{\ell,\beta}.
\end{equation}
No gauge fixing is required to define this finite product of compact-group
integrals.

If $C$ is a closed lattice path based at $y_0$, the ordered product $U(C)$
transforms as $\gamma(y_0)U(C)\gamma(y_0)^{-1}$, by the same cancellation
as Eq.~\eqref{ym:plaquette-gauge}. The real normalized Wilson loop
\begin{equation}\label{ym:wilson-loop-observable}
 W_C(U)=\frac1m\operatorname{Re}\operatorname{tr}U(C)
\end{equation}
is therefore a gauge-invariant observable. Set
$\mathcal O(U)=W_C(U)$ in Eq.~\eqref{ym:observable-expectation}, then insert
Eq.~\eqref{ym:wilson-loop-observable}:
\begin{align}
 \langle W_C\rangle_{\ell,\beta}
 &=\frac{\int W_C(U)e^{-\beta E_\ell(U)}\,\mathcal DU}
        {\int e^{-\beta E_\ell(U)}\,\mathcal DU}\notag\\
 &=\frac{\int \frac1m\operatorname{Re}\operatorname{tr}U(C)
            e^{-\beta E_\ell(U)}\,\mathcal DU}
        {Z_{\ell,\beta}}\notag\\
 &=\frac1{mZ_{\ell,\beta}}
       \int\operatorname{Re}\operatorname{tr}U(C)
             e^{-\beta E_\ell(U)}\,\mathcal DU.
 \label{ym:wilson-loop-expectation}
\end{align}

As an elementary illustration of the integration measure, for $G=U(1)$
write $U(y,\mu)=e^{i\theta_{y,\mu}}$. Then
\begin{align}
 \theta_{y;\mu\nu}
 &=\theta_{y,\mu}+\theta_{y+e_\mu,\nu}
       -\theta_{y+e_\nu,\mu}-\theta_{y,\nu},\\
 1-\operatorname{Re}U_{y;\mu\nu}&=1-\cos\theta_{y;\mu\nu},
 \qquad h(dU)=\frac{d\theta}{2\pi}\quad(0\le\theta<2\pi).
\end{align}
This abelian example explains the measure but is outside the nonabelian
simple-group target of the Millennium problem. For nonabelian groups the
ordered matrix products must be retained.

\subsection{Classical limit of the Wilson energy}
Let $A_\mu\colon\mathbb R^4\to\mathfrak g$ be smooth, $a>0$ the lattice
spacing, and $U(y,\mu)=\exp\big(aA_\mu(ay)\big)$. Write $A_\mu$ for
$A_\mu(ay)$ and $\frac{\partial A_\nu}{\partial x^\mu}$ for its derivative at
$ay$. Taylor expansion gives $A_\nu(ay+ae_\mu)=A_\nu
+a\frac{\partial A_\nu}{\partial x^\mu}+O(a^2)$. With
$e^Xe^Y=\exp\big(X+Y+\frac12[X,Y]+\cdots\big)$, the first two and the last two
factors of Eq.~\eqref{ym:plaquette} are $e^P$ and $e^Q$,
\begin{align}
 P&=a(A_\mu+A_\nu)+a^2\Big(\frac{\partial A_\nu}{\partial x^\mu}
   +\frac12[A_\mu,A_\nu]\Big)+O(a^3),\\
 Q&=-a(A_\mu+A_\nu)+a^2\Big(-\frac{\partial A_\mu}{\partial x^\nu}
   +\frac12[A_\mu,A_\nu]\Big)+O(a^3),
\end{align}
where the commutator in $Q$ is $\frac12[-A_\mu,-A_\nu]$. Then
$[P,Q]=-a^2[A_\mu+A_\nu,A_\mu+A_\nu]+O(a^3)=O(a^3)$ and
\begin{equation}\label{ym:plaquette-expansion}
 U_{y;\mu\nu}=e^Pe^Q=\exp\big(a^2F_{\mu\nu}+O(a^3)\big),\qquad
 F_{\mu\nu}=\frac{\partial A_\nu}{\partial x^\mu}
 -\frac{\partial A_\mu}{\partial x^\nu}+[A_\mu,A_\nu].
\end{equation}
The exponent $M=a^2F_{\mu\nu}+O(a^3)$ lies in $\mathfrak g$, so it is
skew-Hermitian and $\operatorname{tr}M$ is imaginary:
\begin{equation}\label{ym:plaquette-energy}
 1-\frac1m\operatorname{Re}\operatorname{tr}e^M
 =1-\frac1m\operatorname{Re}\operatorname{tr}\Big(1+M+\frac12M^2+O(M^3)\Big)
 =-\frac{a^4}{2m}\operatorname{Re}\operatorname{tr}F_{\mu\nu}^2+O(a^5)
 =\frac{a^4}{2}\|F_{\mu\nu}\|_{\mathfrak g}^2+O(a^5).
\end{equation}
Summing over plaquettes with $a^4\sum_y\to\int d^4x$,
\begin{equation}\label{ym:classical-action}
 \beta E_\ell(U)\ \longrightarrow\ \frac{\beta}{2}\int_{\mathbb R^4}
 \sum_{\mu<\nu}\|F_{\mu\nu}(x)\|_{\mathfrak g}^2\,d^4x
 =\frac{1}{2g_0^2}\int\sum_{\mu<\nu}\|F_{\mu\nu}\|_{\mathfrak g}^2\,d^4x,
 \qquad\beta=\frac{1}{g_0^2}.
\end{equation}
$\beta$ is the inverse squared bare coupling for the inner product
\eqref{ym:inner}. Under the rescaling \eqref{ym:casimir-rescale} the same
action reads $\frac{1}{2g_0'^2}\int\sum\|F\|'^2$ with $g_0'^2=c\,g_0^2$, so
\begin{equation}\label{ym:invariant-combination}
 g_0'^2\,C'=c\,g_0^2\,\frac{C_{\mathfrak g}}{c}=g_0^2\,C_{\mathfrak g}
\end{equation}
does not depend on the normalization of the inner product.

For $G=SU(m)$, choose Hermitian generators with
$\operatorname{tr}(T^aT^b)=\tfrac12\delta^{ab}$ and put $e_a=-iT^a$.
The inner product of Eq.~\eqref{ym:inner} gives
\begin{equation}\label{ym:fundamental-normalization}
 \langle e_a,e_b\rangle_{\mathfrak g}
 =-\frac1m\operatorname{Re}\operatorname{tr}((-iT^a)(-iT^b))
 =\frac{\delta^{ab}}{2m}.
\end{equation}
Thus the component inner product $\langle e_a,e_b\rangle'=\delta^{ab}$
is obtained with $c=2m$. Denoting its coupling by $g=g_0'$, the same action
has
\begin{equation}\label{ym:coupling-dictionary}
 g^2=2m\,g_0^2,\qquad \beta=\frac{1}{g_0^2}=\frac{2m}{g^2}.
\end{equation}
In particular $\beta=4/g^2$ for $SU(2)$ and $\beta=6/g^2$ for $SU(3)$.
Coupling conventions must therefore be compared through the action and
trace normalization, rather than the name of the coupling alone.

\subsection{Asymptotic freedom and the bare coupling}
Choose a basis and adjoint normalization
\begin{equation}\label{ym:adjoint-generator-normalization}
 e_a=-iT^a,\qquad [T^a,T^b]=if^{abc}T^c,\qquad
 f^{acd}f^{bcd}=T(A)\delta^{ab}.
\end{equation}
Equation~\eqref{ym:adjoint-generator-normalization} records the adjoint
convention of Srednicki, Eq.~(70.9) \cite{MSrednicki2007}. Choose the inner product
$\langle e_a,e_b\rangle'=\delta_{ab}$. Then $[e_a,e_b]=-[T^a,T^b]
=-if^{abc}T^c=f^{abc}e_c$, the matrix of $\operatorname{ad}(e_a)$ is
$(\operatorname{ad}e_a)_{cb}=f^{abc}$, and
\begin{equation}\label{ym:killing-structure}
 K(e_a,e_b)=(\operatorname{ad}e_a)_{cd}(\operatorname{ad}e_b)_{dc}
 =f^{adc}f^{bcd}=-f^{acd}f^{bcd}=-T(A)\delta^{ab},\qquad C'=T(A).
\end{equation}
In this normalization Srednicki's Lagrangian
$-\frac14F^{a\mu\nu}F^a_{\mu\nu}$, with the coupling $g$ absorbed into the
field, has Euclidean action
$\frac{1}{4g^2}\int\sum_{\mu,\nu}F^a_{\mu\nu}F^a_{\mu\nu}
=\frac{1}{2g^2}\int\sum_{\mu<\nu}\|F_{\mu\nu}\|'^2$, so $g_0'=g$.
For pure gauge theory, Srednicki, Eq.~(73.41), is
\begin{equation}\label{ym:beta-function}
 \mu\frac{dg}{d\mu}=-b_0'\,g^3+O(g^5),\qquad
 b_0'=\frac{11}{3}\frac{T(A)}{16\pi^2}=\frac{11\,C'}{48\pi^2}.
\end{equation}
Divide by $-\frac12g^3$:
\begin{equation}\label{ym:running-inverse}
 \mu\frac{d}{d\mu}\frac{1}{g^2}=-\frac{2}{g^3}\,\mu\frac{dg}{d\mu}
 =2b_0'+O(g^2),\qquad
 \frac{1}{g^2(\mu)}=\frac{1}{g^2(\mu_0)}+2b_0'\log\frac{\mu}{\mu_0}
 +\cdots .
\end{equation}
The rescaling is $g^2=c\,g_0^2$ and $C'=C_{\mathfrak g}/c$.
Since $c$ is independent of $\mu$, Eq.~\eqref{ym:running-inverse} gives
\begin{equation}\label{ym:running-normalization}
 \mu\frac{d\beta}{d\mu}
 =c\,\mu\frac{d}{d\mu}\frac1{g^2}
 =2c\,b_0'+O(g^2)
 =\frac{11cC'}{24\pi^2}+O(g^2)
 =\frac{11C_{\mathfrak g}}{24\pi^2}+O(g^2).
\end{equation}
Thus in the normalization \eqref{ym:inner}, with $\mu=1/a$ and $\beta=1/g_0^2$,
\begin{equation}\label{ym:slope}
 a\frac{d\beta}{da}=-b_*+\cdots,\qquad
 b_*=2b_0=\frac{11\,C_{\mathfrak g}}{24\pi^2},\qquad
 b_0=\frac{11\,C_{\mathfrak g}}{48\pi^2}.
\end{equation}
For $SU(2)$, $b_*=\frac{11}{3\pi^2}$; for $SU(3)$, $b_*=\frac{33}{4\pi^2}$.

The factor $\frac{1}{8\pi^2}$ per unit of $\log$ is the momentum-shell
integral in four dimensions. The unit sphere $S^3$ has area $2\pi^2$, so for
$L>1$
\begin{equation}\label{ym:shell}
 \int_{1/L\le|k|\le1}\frac{d^4k}{(2\pi)^4}\frac{1}{|k|^4}
 =\frac{2\pi^2}{16\pi^4}\int_{1/L}^{1}\frac{k^3\,dk}{k^4}
 =\frac{\log L}{8\pi^2},\qquad
 b_*=\frac{11\,C_{\mathfrak g}}{3}\cdot\frac{1}{8\pi^2}.
\end{equation}
The angular averages that enter one-loop tensor integrals over $S^3$ follow
from $\langle n_\mu n_\nu\rangle=\frac14\delta_{\mu\nu}$ and
$\langle n_\mu n_\nu n_\rho n_\sigma\rangle=\frac{1}{24}
(\delta_{\mu\nu}\delta_{\rho\sigma}+\delta_{\mu\rho}\delta_{\nu\sigma}
+\delta_{\mu\sigma}\delta_{\nu\rho})$: contracting the second with
$\delta_{\mu\nu}\delta_{\rho\sigma}$ gives $\frac{1}{24}(16+4+4)=1$. In
particular $\langle n_1^2\rangle=\frac14$,
$\langle n_1^2n_2^2\rangle=\frac1{24}$, $\langle n_1^4\rangle=\frac18$.

Integrating Eq.~\eqref{ym:slope}, a continuum limit $a\to0$ at fixed physics
requires
\begin{equation}\label{ym:scaling-law}
 \beta(a)=b_*\log\frac1a+o\Big(\log\frac1a\Big),\qquad
 a\,\Lambda=\exp\Big(-\frac{\beta(a)}{b_*}\Big)\ \text{to leading logarithm},
\end{equation}
where $\Lambda$ is the integration constant, a mass. A physical mass
$m_{\mathrm{phys}}=c\,\Lambda$ corresponds to a lattice mass
\begin{equation}\label{ym:lattice-mass}
 m_{\mathrm{lat}}(\beta)=a\,m_{\mathrm{phys}}\approx c\,e^{-\beta/b_*},
\end{equation}
up to powers of $\beta$ from higher loops: exponentially small in $\beta$.

\subsection{The existence theorem}
\begin{definition}\label{ym:def:family}
A Wilson family for $(m,G,\mathfrak g)$ consists of strictly increasing
sequences of positive integers $(L_p)_{p\in\mathbb N}$ (physical sides) and
$(D_N)_{N\in\mathbb N}$ (inverse spacings, $a_N=1/D_N$), couplings
$\beta_{N,p}\in\mathbb R$, a set $\mathcal T$ of tokens, and for each
$N,p$, $v\in\mathcal T$, $y\in\Lambda_{N,p}=\Lambda_{L_pD_N}$ an insertion
$I_{N,p}(v,y;\cdot)\colon G^{\Lambda_{N,p}\times\{1,\dots,4\}}\to\mathbb C$.
The family is regular if each insertion is continuous, gauge invariant and
depends only on the links in a fixed finite stencil around $y$, and
\begin{equation}\label{ym:reg1}
 \exists\,b_*>0:\qquad
 \lim_{N\to\infty}\ \sup_{p}\Big|\frac{\beta_{N,p}}{\log(1/a_N)}-b_*\Big|=0 .
\end{equation}
\end{definition}
Write $\mu_{N,p}=\mu_{L_pD_N,\beta_{N,p}}$, and $\hat y\in(-L_p/2,L_p/2]^4$
for the physical position of $y\in\Lambda_{N,p}$. For $v\in\mathcal T^n$ and
$f\in C_c^\infty(\mathrm{Conf}_n)$, where $\mathrm{Conf}_n=\{x\in(\mathbb
R^4)^n: x_i\ne x_j\ (i\ne j)\}$, the smeared moment is
\begin{equation}\label{ym:smeared-moment}
 C_{N,p}(v;f)=a_N^{4n}\sum_{y\in\Lambda_{N,p}^n}f(\hat y_1,\dots,\hat y_n)
 \int\prod_{i=1}^nI_{N,p}(v_i,y_i;U)\,\mu_{N,p}(dU).
\end{equation}
\begin{definition}\label{ym:def:ordered}
The family converges along $P\colon\mathbb N\to\mathbb N$ to kernels $K_n(v;x)$
if for every $n,v,f$ there are numbers $R_p$ with
\begin{equation}\label{ym:ordered-limit}
 \lim_{N\to\infty}C_{N,p}(v;f)=R_p\ \ \text{whenever }
 \operatorname{supp}f\subset\big((-\tfrac{L_p}{2},\tfrac{L_p}{2})^4\big)^n,
 \qquad
 \lim_{k\to\infty}R_{P(k)}=\int f(x)\,K_n(v;x)\,dx .
\end{equation}
\end{definition}
The first limit is the ultraviolet limit at fixed volume along the whole
sequence $N=0,1,2,\dots$; the second is the infinite-volume limit along $P$.

\begin{theorem}[Existence; formal proof \cite{OpenAIYM2026}]\label{ym:thm:existence}
Let $(m,G,\mathfrak g)$ be a gauge group with $\mathfrak g$ simple. There is a
regular Wilson family such that for every strictly increasing
$p\colon\mathbb N\to\mathbb N$ there are a strictly increasing
$q\colon\mathbb N\to\mathbb N$ and the following data.
\begin{enumerate}
\item Finite-dimensional spaces $F_d$ of fields of degree $d$, a surjection
 $\pi\colon\mathcal T\to\bigsqcup_dF_d$, a compact metric space $X$ of phases
 with a probability measure $\nu$ of full support and a point $\xi_*\in X$,
 and correlation functions $S^\xi$ ($\xi\in X$) and $S^{\mathrm{mix}}$,
 smooth and permutation symmetric on $\mathrm{Conf}_n$, continuous in $\xi$,
 with
 \begin{equation}\label{ym:mixture}
  S^{\mathrm{mix}}(a;x)=\int_XS^\xi(a;x)\,\nu(d\xi).
 \end{equation}
\item The family converges along $k\mapsto p(q(k))$, in the sense of
 Eq.~\eqref{ym:ordered-limit}, to
 $K_n(v;x)=S^{\mathrm{mix}}(\pi(v_1),\dots,\pi(v_n);x)$.
\item A Hilbert space $\mathcal H$ with vacuum $\Omega$, fields
 $\Phi(\ell,f)$, a unitary representation of the Poincar\'e group and an
 energy--momentum measure supported in the closed forward cone, satisfying
 the Wightman axioms with $\Omega$ the only translation-invariant vector up
 to scalars. The Wightman functions are the boundary values of the analytic
 continuation in the times of $S^{\xi_*}$.
\item An operator product expansion for $S^\xi$ and $S^{\mathrm{mix}}$ with
 coefficient and remainder bounds.
\item A field $\Theta\in F_4$ with $\langle\Omega,\Phi(\Theta,f)\Omega\rangle=0$
 whose two-point function at Euclidean time separation $r$,
 $G_\Theta(r)=S^{\xi_*}(\Theta,\Theta;(0,-\tfrac r2),(0,\tfrac r2))$, satisfies
 \begin{equation}\label{ym:trace-law}
  \lim_{r\to0^+}r^8\Big(\log\frac1r\Big)^2G_\Theta(r)
  =\frac{3\dim_{\mathbb R}\mathfrak g}{\pi^4},
 \end{equation}
 and $G_\Theta$ agrees near $r=0$ with no finite sum of products of
 free-field mass-shell kernels.
\end{enumerate}
In the family of the proof, $b_*$ in Eq.~\eqref{ym:reg1} has the value
\eqref{ym:slope}.
\end{theorem}
The quantifiers matter. The lattice data are chosen before $p$. The
refinement $q$ and all of item~1 and items~3--5 are chosen after $p$ and may
depend on it. The theorem does not assert: that the limit in
Eq.~\eqref{ym:ordered-limit} exists along the full sequence of volumes; that
two subsequences give the same $S^{\mathrm{mix}}$; that $X$ is a point; that
$S^\xi=S^{\xi'}$; or Eq.~\eqref{ym:gap-statement}.

The proof establishes more than the statement
\cite{Armstrong2026,YMGroup2026}: a theory as in item~3 for every phase
$\xi\in X$, with Hamiltonian $H_\xi\ge0$, projection $P_{0,\xi}$ onto
$\mathbb C\Omega_\xi=\ker H_\xi$, and for every finite sum $A$ of products of
smeared fields a constant $M_A$, independent of $\xi$, such that
\begin{equation}\label{ym:low-energy}
 \nu_{A\Omega_\xi}\big((0,u]\big)\le M_A\,u^2\quad(0<u\le1),
\end{equation}
where $\nu_\psi(B)=\langle\psi,E_{H_\xi}(B)\psi\rangle$ is the spectral
measure of $H_\xi$ in the vector $\psi$.

\subsection{Shape of the proof}
The proof has an ultraviolet part and an infinite-volume part
\cite{Armstrong2026}.

\emph{Ultraviolet.} Levels $m=N,N-1,\dots,0$ carry lattices of spacing
$a_m$ on the same torus of side $L_p$; level $N$ is the Wilson lattice and
level $0$ has unit spacing and $L_p^4$ sites. One block-spin step integrates
the fine configuration $U$ against a kernel $k(V,U)\ge0$ that is a
probability density in the coarse configuration $V$:
\begin{equation}\label{ym:block-step}
 \rho_m(V)=\int k(V,U)\,\rho_{m+1}(U)\,dU,\qquad\int k(V,U)\,dV=1 .
\end{equation}
By Tonelli's theorem,
\begin{equation}\label{ym:partition-preserved}
 \int\rho_m(V)\,dV=\int\Big(\int k(V,U)\,dV\Big)\rho_{m+1}(U)\,dU
 =\int\rho_{m+1}(U)\,dU,\qquad
 \int\rho_0\,dV=\int\rho_N\,dU=Z_{L_pD_N,\beta_{N,p}} .
\end{equation}
The density at level $m$ is controlled in the form
\begin{equation}\label{ym:effective-density}
 \rho_m(U)=\exp\big(v|\Lambda_m|-\beta_{N,p,m}\,\mathcal K(U)-\mathcal L(U)
 -\mathcal E(U)\big)\,\mathcal G(U),
\end{equation}
with $\mathcal K$ a gauge-invariant local action normalized so that its
quadratic part at zero momentum is that of Eq.~\eqref{ym:classical-action},
$\mathcal L$ the renormalized insertions, $\mathcal E$ irrelevant terms and
$\mathcal G$ a large-field factor. The bare coupling $\beta_{N,p}$ is tuned
so that the level-$0$ coupling is one fixed large number $\bar\beta$:
\begin{equation}\label{ym:running}
 \beta_{N,p,0}=\bar\beta,\qquad
 \beta_{N,p,m}=\bar\beta+b_*\log\frac{1}{a_m}+o\Big(\log\frac{1}{a_m}\Big)
 \quad\text{uniformly in }N\ge m\text{ and }p .
\end{equation}
At $m=N$ this is Eq.~\eqref{ym:reg1}. Bounds on moments of insertions that
are uniform in $N$ and $p$ give the first limit in
Eq.~\eqref{ym:ordered-limit}.

\emph{Infinite volume.} The uniform bounds and their difference quotients
give equicontinuity on compact subsets of $\mathrm{Conf}_n$; a diagonal
Arzel\`a--Ascoli argument extracts $q$. Reflection positivity passes to the
limit. The limit $S^{\mathrm{mix}}$ need not cluster; it is decomposed over
the spectrum $X$ of a commuting family of operators as in
Eq.~\eqref{ym:mixture}, and each $S^\xi$ is reconstructed.

No estimate in the infinite-volume part depends on the dynamics at distances
beyond the unit scale of level $0$. By Eqs.~\eqref{ym:lattice-mass} and
\eqref{ym:running}, the level-$0$ system is a lattice gauge theory at
coupling $\bar\beta$ whose expected mass in lattice units is of order
$e^{-\bar\beta/b_*}$.

\subsection{Spectral measures and the mass gap}
Let $H\ge0$ be self-adjoint on $\mathcal H$ with spectral projections
$E_H(B)$, $P_0=E_H(\{0\})$, and for $\psi\in\mathcal H$ let
$\nu_\psi(B)=\langle\psi,E_H(B)\psi\rangle$ and
$N_\psi(u)=\nu_\psi((0,u])$. For $s>0$,
\begin{equation}\label{ym:laplace}
 \big\langle\psi,(e^{-sH}-P_0)\psi\big\rangle
 =\int_{(0,\infty)}e^{-sE}\,\nu_\psi(dE).
\end{equation}
\begin{lemma}\label{ym:lem:layer}
$\displaystyle\int_{(0,\infty)}e^{-sE}\,\nu_\psi(dE)
=\int_0^\infty s\,e^{-su}\,N_\psi(u)\,du$.
\end{lemma}
\begin{proof}
$e^{-sE}=\int_E^\infty s\,e^{-su}\,du$. Insert this and exchange the
integrals by Tonelli's theorem:
\begin{equation}
 \int_{(0,\infty)}\int_0^\infty\mathbf 1_{\{E\le u\}}\,s\,e^{-su}\,du\,
 \nu_\psi(dE)
 =\int_0^\infty s\,e^{-su}\,\nu_\psi\big((0,u]\big)\,du .
\end{equation}
\end{proof}
\begin{proposition}[Polynomial decay from Eq.~\eqref{ym:low-energy}]
\label{ym:prop:inverse-square}
If $N_\psi(u)\le Du^2$ for all $u>0$, then
$0\le\langle\psi,(e^{-sH}-P_0)\psi\rangle\le2D\,s^{-2}$ for all $s>0$.
\end{proposition}
\begin{proof}
By Eq.~\eqref{ym:laplace} and Lemma~\ref{ym:lem:layer},
\begin{equation}
 \big\langle\psi,(e^{-sH}-P_0)\psi\big\rangle
 \le D\,s\int_0^\infty u^2e^{-su}\,du
 =D\,s\,\frac{\Gamma(3)}{s^3}=\frac{2D}{s^2}.
\end{equation}
\end{proof}
The hypothesis follows from Eq.~\eqref{ym:low-energy} with
$D=\max(M_A,\|A\Omega_\xi\|^2)$: for $u>1$,
$N_\psi(u)\le\|\psi\|^2\le\|\psi\|^2u^2$.

\begin{proposition}[Gap and exponential decay]\label{ym:prop:gap-decay}
Let $\Delta>0$ and let $\mathcal D\subset\mathcal H$ be dense. The following
are equivalent.
\begin{enumerate}
\item $\operatorname{spec}(H)\subset\{0\}\cup[\Delta,\infty)$.
\item For every $\psi\in\mathcal D$ there are $C_\psi$ and $s_\psi$ with
 $\langle\psi,(e^{-sH}-P_0)\psi\rangle\le C_\psi e^{-s\Delta}$ for all
 $s\ge s_\psi$.
\end{enumerate}
\end{proposition}
\begin{proof}
(1)$\Rightarrow$(2): $\nu_\psi((0,\Delta))=0$, so Eq.~\eqref{ym:laplace} is
at most $e^{-s\Delta}\nu_\psi([\Delta,\infty))\le e^{-s\Delta}\|\psi\|^2$.

(2)$\Rightarrow$(1): let $0<\Delta'<\Delta$ and $\psi\in\mathcal D$. On
$(0,\Delta']$, $e^{-sE}\ge e^{-s\Delta'}$, so for $s\ge s_\psi$
\begin{equation}
 e^{-s\Delta'}\,\nu_\psi\big((0,\Delta']\big)
 \le\int_{(0,\infty)}e^{-sE}\,\nu_\psi(dE)\le C_\psi e^{-s\Delta},
 \qquad
 \nu_\psi\big((0,\Delta']\big)\le C_\psi e^{-s(\Delta-\Delta')}
 \xrightarrow[s\to\infty]{}0 .
\end{equation}
Hence $\|E_H((0,\Delta'])\psi\|^2=0$ on the dense set $\mathcal D$; the
projection is bounded, so $E_H((0,\Delta'])=0$ for every $\Delta'<\Delta$,
and $E_H((0,\Delta))=0$.
\end{proof}
For $\psi=A\Omega$ with $A$ a product of fields smeared at times $\le0$,
$\langle\psi,(e^{-sH}-P_0)\psi\rangle$ is a truncated Euclidean correlation
function at time separation $s$. So Eq.~\eqref{ym:gap-statement} is
exponential decay of truncated correlations of $S^\xi$, with one rate for
all local fields.

Eq.~\eqref{ym:low-energy} does not decide the question. A massless particle
saturates it: for a one-particle state of a massless scalar field with
momentum-space wave function $\hat f$, energy $E=|\mathbf p|$ and invariant
measure $\frac{d^3p}{2|\mathbf p|}$,
\begin{equation}\label{ym:massless-rate}
 N_\psi(u)=\int_{|\mathbf p|\le u}|\hat f(\mathbf p)|^2
 \frac{d^3p}{2|\mathbf p|}
 \le\sup|\hat f|^2\cdot4\pi\int_0^u\frac{p^2\,dp}{2p}
 =\pi\sup|\hat f|^2\,u^2 ,
\end{equation}
with equality in the limit $u\to0$ when $\hat f$ is continuous and
$\hat f(0)\ne0$. A gapped theory has $N_\psi(u)=0$ for $u<\Delta$. Both
satisfy Eq.~\eqref{ym:low-energy}.

\subsection{Transfer operator and the purity of the trace}
Consider the torus with sides $(n_1,n_2,n_3,n_4)$, the same $\beta\ge0$, and
partition function $Z(n_1,n_2,n_3,n_4)$ as in
Eq.~\eqref{ym:wilson-measure}. $E_\ell$ sums over all six planes, so
relabelling the axes shows that $Z$ is symmetric in its four arguments.

Split a configuration into time slices: spatial links
$U_t=(U((\mathbf x,t),j))_{\mathbf x,j\le3}$ and temporal links
$V_t(\mathbf x)=U((\mathbf x,t),4)$. Let $E_{\mathrm s}(U_t)$ be the sum of
Eq.~\eqref{ym:wilson-energy} over the plaquettes in the planes $(j,k)$,
$j<k\le3$, of slice $t$. By Eq.~\eqref{ym:plaquette} the plaquette in the
plane $(j,4)$ at $(\mathbf x,t)$ is
$U_t(\mathbf x,j)\,V_t(\mathbf x+e_j)\,U_{t+1}(\mathbf x,j)^{-1}\,V_t(\mathbf
x)^{-1}$. Define the kernel
\begin{align}
 \mathsf T(U,U')&=e^{-\frac\beta2E_{\mathrm s}(U)}\,\mathsf K(U,U')\,
 e^{-\frac\beta2E_{\mathrm s}(U')},\label{ym:transfer-kernel}\\
 \mathsf K(U,U')&=\int\prod_{\mathbf x}h\big(dV(\mathbf x)\big)
 \exp\Big(-\beta\sum_{\mathbf x,j}\Big(1-\frac1m\operatorname{Re}
 \operatorname{tr}\,U(\mathbf x,j)V(\mathbf x+e_j)U'(\mathbf x,j)^{-1}
 V(\mathbf x)^{-1}\Big)\Big).\notag
\end{align}
Each spatial plaquette of slice $t$ contributes
$\frac\beta2+\frac\beta2$ through the two kernels adjacent to $U_t$, so
\begin{equation}\label{ym:trace-formula}
 Z(n_1,n_2,n_3,n)=\int\prod_{t\in\mathbb Z/n}dU_t\
 \mathsf T(U_0,U_1)\mathsf T(U_1,U_2)\cdots\mathsf T(U_{n-1},U_0)
 =\operatorname{Tr}\mathsf T^{\,n},
\end{equation}
where $\mathsf T$ acts on $L^2$ of the spatial configurations on the torus
$(n_1,n_2,n_3)$. The kernel is real, continuous on a compact space, and
symmetric. For the temporal plaquette
$W=U(\mathbf x,j)V(\mathbf x+e_j)U'(\mathbf x,j)^{-1}V(\mathbf x)^{-1}$, which
is unitary,
\begin{equation}\label{ym:kernel-symmetry}
 \operatorname{Re}\operatorname{tr}W=\operatorname{Re}\operatorname{tr}W^{-1}
 =\operatorname{Re}\operatorname{tr}\,U'(\mathbf x,j)V(\mathbf x+e_j)^{-1}
 U(\mathbf x,j)^{-1}V(\mathbf x),
\end{equation}
by cyclicity of the trace. The right side is the temporal plaquette of
$\mathsf K(U',U)$ with $V$ replaced by $V^{-1}$, and Haar measure on a
compact group is invariant under inversion; so
$\mathsf K(U,U')=\mathsf K(U',U)$ and $\mathsf T$ is self-adjoint. It is
nonnegative and trace class by Proposition~\ref{ym:sc-prop:positivity}, whose
proof does not use the trace formula; in particular $\operatorname{Tr}\mathsf T$ in
Eq.~\eqref{ym:trace-formula} with $n=1$ is the integral of the diagonal.
\begin{proposition}[Purity]\label{ym:prop:purity}
Let $\mathsf T$ be a self-adjoint trace-class operator with eigenvalues
$\lambda_1\ge\lambda_2\ge\dots\ge0$, $\lambda_1>0$, and $Z(n)=\operatorname{Tr}
\mathsf T^{\,n}$. Put
\begin{equation}\label{ym:purity-defs}
 r_j=\Big(\frac{\lambda_j}{\lambda_1}\Big)^n\in[0,1],\qquad
 s_n=\sum_{j\ge2}r_j,\qquad q_n=\sum_{j\ge2}r_j^2,\qquad
 \mathsf P(n)=\frac{Z(2n)}{Z(n)^2}.
\end{equation}
Then
\begin{equation}\label{ym:purity-bounds}
 \mathsf P(n)=\frac{1+q_n}{(1+s_n)^2}\le1,\qquad
 s_n\ \le\ \frac{1}{\mathsf P(n)}-1\ \le\ 2s_n+s_n^2,\qquad
 s_{2n}\le s_n^2 .
\end{equation}
Consequently the ratio of the two largest eigenvalues satisfies
\begin{equation}\label{ym:gap-from-purity}
 \Big(\frac{\lambda_2}{\lambda_1}\Big)^n=r_2\le s_n\le\frac{1}{\mathsf P(n)}-1 .
\end{equation}
\end{proposition}
\begin{proof}
$Z(n)=\sum_j\lambda_j^n=\lambda_1^n(1+s_n)$ and
$Z(2n)=\lambda_1^{2n}(1+q_n)$ give the formula for $\mathsf P(n)$. Since
$0\le r_j\le1$, $r_j^2\le r_j$ and $0\le q_n\le s_n$; hence
$\mathsf P(n)\le\frac{1+s_n}{(1+s_n)^2}\le1$. Next,
\begin{equation}
 \frac{1}{\mathsf P(n)}-1-s_n
 =\frac{(1+s_n)^2-(1+q_n)-s_n(1+q_n)}{1+q_n}
 =\frac{s_n+s_n^2-q_n-s_nq_n}{1+q_n}
 =\frac{(1+s_n)(s_n-q_n)}{1+q_n}\ \ge0,
\end{equation}
and $\frac{1}{\mathsf P(n)}-1=\frac{2s_n+s_n^2-q_n}{1+q_n}\le2s_n+s_n^2$.
Finally $s_{2n}=\sum_{j\ge2}r_j^2=q_n\le\big(\sum_{j\ge2}r_j\big)^2=s_n^2$.
For Eq.~\eqref{ym:gap-from-purity}, $r_2$ is one of the nonnegative terms of
$s_n$, so $r_2\le s_n$, and $s_n\le\mathsf P(n)^{-1}-1$ was shown above.
\end{proof}
The transfer gap is
\begin{equation}\label{ym:transfer-gap}
 \mathsf m=\log\frac{\lambda_1}{\lambda_2}\quad(\lambda_2>0),\qquad
 \mathsf m=+\infty\quad(\lambda_2=0).
\end{equation}
For $\lambda_2>0$ the right side of Eq.~\eqref{ym:gap-from-purity} is at least
$r_2>0$, so both sides have logarithms:
\begin{equation}\label{ym:gap-from-purity-log}
 -n\,\mathsf m=n\log\frac{\lambda_2}{\lambda_1}\le
 \log\Big(\frac{1}{\mathsf P(n)}-1\Big),\qquad
 \mathsf m\ \ge\ \frac1n\log\frac{1}{\mathsf P(n)^{-1}-1},
\end{equation}
and the bound is positive when $\mathsf P(n)>\frac12$. The ratio form
Eq.~\eqref{ym:gap-from-purity} also holds when $\lambda_2=0$, where the
logarithm of the ratio does not exist. This happens at $\beta=0$: every
exponential in Eq.~\eqref{ym:transfer-kernel} equals $1$, the Haar measures
are normalized, and
\begin{equation}\label{ym:rank-one}
 \mathsf K(U,U')=\int\prod_{\mathbf x}h\big(dV(\mathbf x)\big)\cdot1=1,\qquad
 \mathsf T(U,U')=1,\qquad
 (\mathsf Tf)(U)=\int f(U')\,dU' .
\end{equation}
$\mathsf T$ is the projection onto constant functions: $\lambda_1=1$ and
$\lambda_j=0$ for $j\ge2$. Then $s_n=q_n=0$, $\mathsf P(n)=1$ and both sides
of Eq.~\eqref{ym:gap-from-purity} vanish, while $\mathsf m=+\infty$.
The sum $s_n$ counts every excited state with its multiplicity, in every
symmetry sector. At fixed $n$ it grows with the spatial volume. The defect of
a hypercube controls it on the doubled spatial torus.
\begin{definition}\label{ym:def:defect}
The hypercubic defect at side $n$ is
\begin{equation}\label{ym:defect}
 \Delta_\beta(n)=16\log Z(n,n,n,n)-\log Z(2n,2n,2n,2n).
\end{equation}
\end{definition}
For comparison, the square-defect expression of the two-dimensional
model in \cite{OpenAIO4Sharp2026}, Eq.~(1.6), is
\begin{equation}\label{ym:two-dimensional-square-defect}
 4\log Z(n,n)-\log Z(2n,2n).
\end{equation}
Equation~\eqref{ym:defect} is the four-dimensional counterpart of
Eq.~\eqref{ym:two-dimensional-square-defect}: a doubled square has four
original squares, whereas a doubled four-dimensional hypercube has
$2^4=16$ original hypercubes.
\begin{proposition}[Defect and gap on the doubled torus]\label{ym:prop:defect}
On every spatial torus the operator $\mathsf T$ of
Eq.~\eqref{ym:transfer-kernel} is nonnegative
(Proposition~\ref{ym:sc-prop:positivity}). Let
\begin{align}
 \mathsf P_1&=\frac{Z(2n,n,n,n)}{Z(n,n,n,n)^2},&
 \mathsf P_2&=\frac{Z(2n,2n,n,n)}{Z(2n,n,n,n)^2},\notag\\
 \mathsf P_3&=\frac{Z(2n,2n,2n,n)}{Z(2n,2n,n,n)^2},&
 \mathsf P_4&=\frac{Z(2n,2n,2n,2n)}{Z(2n,2n,2n,n)^2}.\label{ym:four-purities}
\end{align}
Then $0<\mathsf P_k\le1$,
\begin{equation}\label{ym:defect-telescope}
 e^{-\Delta_\beta(n)}=\mathsf P_1^{\,8}\,\mathsf P_2^{\,4}\,\mathsf P_3^{\,2}\,
 \mathsf P_4,\qquad\Delta_\beta(n)\ge0,\qquad
 \mathsf P_4\ge e^{-\Delta_\beta(n)},
\end{equation}
and the two largest eigenvalues $\lambda_1(2n)\ge\lambda_2(2n)$ of $\mathsf T$
on the spatial torus $(2n,2n,2n)$ satisfy
\begin{equation}\label{ym:gap-from-defect}
 \Big(\frac{\lambda_2(2n)}{\lambda_1(2n)}\Big)^n\ \le\ e^{\Delta_\beta(n)}-1 .
\end{equation}
\end{proposition}
\begin{proof}
By the symmetry of $Z$ in its arguments, each $\mathsf P_k$ is a purity
$Z(2n)/Z(n)^2$ of Eq.~\eqref{ym:purity-defs} for the transfer operator in the
direction being doubled, on the torus formed by the other three sides; so
$0<\mathsf P_k\le1$ by Eq.~\eqref{ym:purity-bounds}. Multiply:
\begin{equation}
 \mathsf P_1^{\,8}\mathsf P_2^{\,4}\mathsf P_3^{\,2}\mathsf P_4
 =\frac{Z(2n,n,n,n)^8}{Z(n,n,n,n)^{16}}\cdot
  \frac{Z(2n,2n,n,n)^4}{Z(2n,n,n,n)^{8}}\cdot
  \frac{Z(2n,2n,2n,n)^2}{Z(2n,2n,n,n)^{4}}\cdot
  \frac{Z(2n,2n,2n,2n)}{Z(2n,2n,2n,n)^{2}}
 =\frac{Z(2n,2n,2n,2n)}{Z(n,n,n,n)^{16}} .
\end{equation}
All factors are at most $1$, so $\Delta_\beta(n)\ge0$ and
$\mathsf P_4\ge\mathsf P_1^{\,8}\mathsf P_2^{\,4}\mathsf P_3^{\,2}\mathsf P_4
=e^{-\Delta_\beta(n)}$. $\mathsf P_4$ is the purity at time extent $n$ on the
spatial torus $(2n,2n,2n)$, and $\mathsf P_4\ge e^{-\Delta_\beta(n)}>0$ gives
$\mathsf P_4^{-1}-1\le e^{\Delta_\beta(n)}-1$. Eq.~\eqref{ym:gap-from-purity}
for this purity then gives Eq.~\eqref{ym:gap-from-defect}.
\end{proof}
When $\lambda_2(2n)>0$, taking logarithms in Eq.~\eqref{ym:gap-from-defect}
as in Eq.~\eqref{ym:gap-from-purity-log} gives
\begin{equation}\label{ym:gap-from-defect-log}
 \mathsf m(2n)\ \ge\ \frac1n\log\frac{1}{e^{\Delta_\beta(n)}-1},
\end{equation}
which is positive when $\Delta_\beta(n)<\log2$. A vanishing defect forces
$\lambda_2(2n)=0$: Eq.~\eqref{ym:gap-from-defect} reads
$(\lambda_2/\lambda_1)^n\le e^0-1=0$.
A bulk free energy $f\,n_1n_2n_3n_4$ in $\log Z$ cancels from
Eq.~\eqref{ym:defect}: $16n^4-(2n)^4=0$. The defect measures only the
deviation of $\log Z$ from extensivity.

\subsection{What a proof of the gap requires}
In two dimensions, \cite{OpenAIO4Sharp2026} states for the $O(4)$ model that
the square defect of Eq.~\eqref{ym:two-dimensional-square-defect} obeys
$\Delta_\beta(2n)\le64\,\Delta_\beta(n)^2$ in the small-defect regime, and that
for even $n\ge4$, $\Delta_\beta(n)\le2^{-10}$ implies
$m_{\mathrm{lat}}(\beta)\ge\frac{\log2}{n}$ on the dyadic circumferences
$2^kn$ and their local limits; a separate estimate makes the defect small on
a box of bounded physical size, and the renormalization map carries the test
to all cutoffs. For Wilson lattice gauge theory the first statement holds in
the four-dimensional form of Theorem~\ref{ym:thm:recurrence}; a small defect
on a box of bounded physical size is open.

\paragraph{Doubling costs.} For sides $x=(x_1,x_2,x_3,x_4)$ with
$x_i\ge1$, write $F(x)=\log Z(x_1,x_2,x_3,x_4)$ and let $D_ix$ be $x$ with
side $i$ doubled. The cost of doubling side $i$ is
\begin{equation}\label{ym:doubling-cost}
 p_i(x)=2F(x)-F(D_ix)=-\log\frac{Z(D_ix)}{Z(x)^2}.
\end{equation}
By the symmetry of $Z$, the ratio $Z(D_ix)/Z(x)^2$ is the purity
$\mathsf P(t)=Z(2t)/Z(t)^2$ of Eq.~\eqref{ym:purity-defs} at time extent
$t=x_i$ for the transfer operator in direction $i$ on the torus formed by the
other three sides. That operator is nonnegative and trace class
(Proposition~\ref{ym:sc-prop:positivity}), so Eq.~\eqref{ym:purity-bounds}
applies and
\begin{equation}\label{ym:cost-nonneg}
 p_i(x)=-\log\mathsf P(x_i)\ \ge\ 0 .
\end{equation}
\begin{lemma}[One doubling]\label{ym:lem:one-doubling}
If $p_i(x)\le\frac14$, then $p_i(D_ix)\le4\,p_i(x)^2$.
\end{lemma}
\begin{proof}
Fix the transverse torus and write $t=x_i$, $p=p_i(x)$, $s=s_t$. By
Eq.~\eqref{ym:purity-bounds},
\begin{equation}
 s\ \le\ \frac{1}{\mathsf P(t)}-1=e^{p}-1=\int_0^pe^u\,du\ \le\ p\,e^{p}.
\end{equation}
$D_ix$ has the same transverse torus and time extent $2t$, so
$p_i(D_ix)=-\log\mathsf P(2t)$. Eq.~\eqref{ym:purity-bounds} at $2t$, and
$s_{2t}\le s^2$ from Eq.~\eqref{ym:purity-bounds} at $t$, give
\begin{align}
 p_i(D_ix)=\log\frac{1}{\mathsf P(2t)}
 &\le\log\big(1+2s_{2t}+s_{2t}^2\big)=2\log(1+s_{2t})\le2s_{2t}\notag\\
 &\le2s^2\le2p^2e^{2p}\le4p^2,
\end{align}
the last step because $e^{2p}\le e^{1/2}<2$ for $p\le\frac14$.
\end{proof}
\begin{lemma}[Commuting doublings]\label{ym:lem:commuting}
For $i\ne j$,
\begin{equation}\label{ym:commuting-doublings}
 p_i(D_jx)=2p_i(x)-2p_j(x)+p_j(D_ix).
\end{equation}
\end{lemma}
\begin{proof}
$D_iD_jx=D_jD_ix$. Insert Eq.~\eqref{ym:doubling-cost} on the right:
\begin{align}
 &2\big(2F(x)-F(D_ix)\big)-2\big(2F(x)-F(D_jx)\big)
 +\big(2F(D_ix)-F(D_jD_ix)\big)\notag\\
 &\quad=4F(x)-2F(D_ix)-4F(x)+2F(D_jx)+2F(D_ix)-F(D_iD_jx)\notag\\
 &\quad=2F(D_jx)-F(D_i(D_jx))=p_i(D_jx).
\end{align}
\end{proof}
\paragraph{The defect as a path sum.} Double the four sides of
$x^0=(n,n,n,n)$ one at a time: $x^1=D_1x^0$, $x^2=D_2x^1$, $x^3=D_3x^2$,
$x^4=D_4x^3=(2n,2n,2n,2n)$. With Eq.~\eqref{ym:doubling-cost},
\begin{align}
 &8p_1(x^0)+4p_2(x^1)+2p_3(x^2)+p_4(x^3)\notag\\
 &\quad=16F(x^0)-8F(x^1)+8F(x^1)-4F(x^2)+4F(x^2)-2F(x^3)+2F(x^3)-F(x^4)
 \notag\\
 &\quad=16F(x^0)-F(x^4)=\Delta_\beta(n),\label{ym:defect-path}
\end{align}
and the same holds for the four directions taken in any order. Every term is
nonnegative by Eq.~\eqref{ym:cost-nonneg} and every weight is at least $1$,
so each doubling cost on such a path is at most $\Delta_\beta(n)$. A pair
$(i,x)$ with $x_i=n$ and $x_j\in\{n,2n\}$ lies on such a path: first double
the sides $j$ with $x_j=2n$, then side $i$, then the rest. Hence
\begin{equation}\label{ym:cost-cube}
 p_i(x)\le\Delta_\beta(n)\qquad\text{when }x_i=n\text{ and }x_j\in\{n,2n\}
 \text{ for }j\ne i.
\end{equation}
\begin{theorem}[Defect recurrence]\label{ym:thm:recurrence}
For every compact $G\subset U(m)$, $\beta\ge0$ and $n\ge1$,
\begin{equation}\label{ym:recurrence}
 \Delta_\beta(n)\le\eta_*:=\frac{1}{7480}\quad\Longrightarrow\quad
 \Delta_\beta(2n)\le c_1\,\Delta_\beta(n)^2,\qquad c_1=3740=\frac{1}{2\eta_*}.
\end{equation}
\end{theorem}
The argument extends the two-direction rectangle identity of
\cite{OpenAIO4Sharp2026} to the four directions of the hypercube; it was
first written out in an independent review of these notes.
\begin{proof}
Write $\delta=\Delta_\beta(n)$. We show: if $x_i=n$, every other side is $n$,
$2n$ or $4n$, and exactly $j$ of them equal $4n$ ($0\le j\le3$), then
\begin{equation}\label{ym:cost-induction}
 p_i(x)\le A_j:=3^j\delta .
\end{equation}
For $j=0$ this is Eq.~\eqref{ym:cost-cube}. For $j\ge1$ choose $h$ with
$x_h=4n$ and let $y$ be $x$ with side $h$ halved, so $x=D_hy$ and $y_h=2n$.
Eq.~\eqref{ym:commuting-doublings} for the pair $(i,h)$ at $y$, with
$p_h(y)\ge0$, gives
\begin{equation}
 p_i(x)=p_i(D_hy)=2p_i(y)-2p_h(y)+p_h(D_iy)\le2p_i(y)+p_h(D_iy).
\end{equation}
$y$ has $y_i=n$ and $j-1$ sides of length $4n$, so $p_i(y)\le A_{j-1}$. Side
$h$ of $D_iy$ has length $2n$; let $z$ be $D_iy$ with side $h$ halved, so
$D_iy=D_hz$, $z_h=n$, side $i$ of $z$ is $2n$, and $z$ has $j-1$ sides of
length $4n$. Eq.~\eqref{ym:cost-induction} in direction $h$ gives
$p_h(z)\le A_{j-1}$, and $A_{j-1}\le9\delta\le\frac14$, so by
Lemma~\ref{ym:lem:one-doubling}
\begin{equation}
 p_h(D_iy)=p_h(D_hz)\le4p_h(z)^2\le4A_{j-1}^2\le A_{j-1}.
\end{equation}
Together, $p_i(x)\le2A_{j-1}+A_{j-1}=3A_{j-1}=A_j$.

Now apply Eq.~\eqref{ym:defect-path} at side $2n$:
$\Delta_\beta(2n)=8p_1(y^0)+4p_2(y^1)+2p_3(y^2)+p_4(y^3)$ with
\begin{equation}
 y^0=(2n,2n,2n,2n),\quad y^1=(4n,2n,2n,2n),\quad y^2=(4n,4n,2n,2n),\quad
 y^3=(4n,4n,4n,2n).
\end{equation}
Side $k$ of $y^{k-1}$ has length $2n$, so $y^{k-1}=D_kw^{k-1}$ with
\begin{equation}
 w^0=(n,2n,2n,2n),\quad w^1=(4n,n,2n,2n),\quad w^2=(4n,4n,n,2n),\quad
 w^3=(4n,4n,4n,n),
\end{equation}
which have $0,1,2,3$ sides of length $4n$. By Eq.~\eqref{ym:cost-induction},
$p_k(w^{k-1})\le A_{k-1}\le27\delta\le\frac14$, since
$27\eta_*=\frac{27}{7480}<\frac14$. Lemma~\ref{ym:lem:one-doubling} gives
$p_k(y^{k-1})=p_k(D_kw^{k-1})\le4A_{k-1}^2$, and
\begin{align}
 \Delta_\beta(2n)&\le8\cdot4A_0^2+4\cdot4A_1^2+2\cdot4A_2^2+4A_3^2
 =4\delta^2\big(8+4\cdot9+2\cdot81+729\big)\notag\\
 &=4\cdot935\,\delta^2=3740\,\delta^2 .
\end{align}
\end{proof}
\begin{conjecture}[Small defect at a bounded physical scale]\label{ym:conj:defect}
There is $\ell_*<\infty$ such that along the couplings of
Eq.~\eqref{ym:scaling-law} there are integers $n(\beta)\le\ell_*\,
e^{\beta/b_*}$ with $\Delta_\beta(n(\beta))\le\eta_*$, the same bare
parameters entering both boxes of Eq.~\eqref{ym:defect}.
\end{conjecture}
\begin{proposition}[Conditional lattice gap]\label{ym:prop:conditional-gap}
If $\Delta_\beta(n)\le\eta_*$ for some $n\ge1$, then for every $k\ge0$ the two
largest eigenvalues of $\mathsf T$ on the spatial torus of side $2^{k+1}n$
satisfy
\begin{equation}\label{ym:conditional-gap}
 \Big(\frac{\lambda_2(2^{k+1}n)}{\lambda_1(2^{k+1}n)}\Big)^n\le\frac12,
 \qquad\text{that is}\qquad
 \mathsf m(2^{k+1}n)\ \ge\ \frac{\log2}{n}.
\end{equation}
\end{proposition}
\begin{proof}
Put $\delta_k=\Delta_\beta(2^kn)\ge0$. Then $c_1\delta_0\le\frac12$, and if
$c_1\delta_k\le\frac12$ then $\delta_k\le\eta_*$ and
Eq.~\eqref{ym:recurrence} gives
\begin{equation}
 c_1\delta_{k+1}\le c_1\cdot c_1\delta_k^2=(c_1\delta_k)^2\le\tfrac14 .
\end{equation}
By induction $c_1\delta_k\le(c_1\delta_0)^{2^k}\le2^{-2^k}$ for all $k$.
For $0\le x\le1$, convexity of $e^x$ gives $e^x-1\le(e-1)x\le2x$. Apply
Eq.~\eqref{ym:gap-from-defect} at side $2^kn$, with $r=\lambda_2/\lambda_1$
on the spatial torus of side $2^{k+1}n$:
\begin{equation}
 \big(r^n\big)^{2^k}=r^{2^kn}\le e^{\delta_k}-1\le2\delta_k
 \le\frac{2}{c_1}\,2^{-2^k}\le\Big(\frac12\Big)^{2^k},
\end{equation}
since $c_1\ge2$. Both sides are nonnegative, so taking the $2^k$-th root gives
$r^n\le\frac12$. If $\lambda_2>0$ this is $\mathsf m\ge\frac{\log2}{n}$ by
Eq.~\eqref{ym:transfer-gap}; if $\lambda_2=0$, $\mathsf m=+\infty$.
\end{proof}
The double-exponential decay of $\delta_k$ compensates the factor
$\frac{1}{2^kn}$, and no positivity of the defects is needed.

\paragraph{A seed at strong coupling.}
\begin{proposition}[Strong-coupling seed]\label{ym:prop:strong-seed}
For every compact $G\subset U(m)$, $\beta\ge0$ and $n\ge1$,
\begin{equation}\label{ym:strong-seed}
 \Delta_\beta(n)\le192\,\beta\,n^4 .
\end{equation}
Hence $\Delta_\beta(n)\le\eta_*$ for $0\le\beta\le\frac{1}{7480\cdot192\,n^4}$,
and Eq.~\eqref{ym:conditional-gap} holds at those couplings.
\end{proposition}
\begin{proof}
On the torus with sides $x$ there is one plaquette for each site and each of
the $\binom42=6$ planes, so the action has
$P(x)=6\,x_1x_2x_3x_4$ terms. For unitary $U_p$,
$|\operatorname{Re}\operatorname{tr}U_p|\le m$, so each term satisfies
\begin{equation}
 0\ \le\ 1-\frac1m\operatorname{Re}\operatorname{tr}U_p\ \le\ 2,
 \qquad
 e^{-2\beta P(x)}\ \le\ \exp\Big(-\beta\sum_p\Big(1-\frac1m
 \operatorname{Re}\operatorname{tr}U_p\Big)\Big)\ \le\ 1 .
\end{equation}
Integrating against the normalized product Haar measure of
Eq.~\eqref{ym:wilson-measure},
\begin{equation}\label{ym:log-z-bounds}
 -2\beta P(x)\ \le\ \log Z(x)\ \le\ 0 .
\end{equation}
In Eq.~\eqref{ym:defect}, bound $\log Z(n,n,n,n)\le0$ from above and
$\log Z(2n,2n,2n,2n)\ge-2\beta P(2n,2n,2n,2n)$ from below:
\begin{equation}
 \Delta_\beta(n)\le16\cdot0+2\beta\cdot6\,(2n)^4=192\,\beta\,n^4 .
\end{equation}
For $\beta\le\frac{1}{7480\cdot192\,n^4}$ the right side is at most
$\frac1{7480}=\eta_*$, and Proposition~\ref{ym:prop:conditional-gap} applies.
\end{proof}
This is the regime in which convergent strong-coupling expansions already give
a mass gap \cite{OsterwalderSeiler1978}. Proposition~\ref{ym:prop:strong-seed}
shows that the chain of Propositions~\ref{ym:prop:purity},
\ref{ym:prop:defect}, \ref{ym:prop:conditional-gap} and
Theorem~\ref{ym:thm:recurrence} closes there with no further input; it says
nothing about boxes of bounded physical size as $\beta\to\infty$, which is
Conjecture~\ref{ym:conj:defect}. With Conjecture~\ref{ym:conj:defect}, in
physical units,
\begin{equation}\label{ym:physical-gap}
 \frac{\mathsf m}{a}\ \ge\ \frac{\log2}{a\,n(\beta)}
 \ \ge\ \frac{\log2}{\ell_*\,a\,e^{\beta/b_*}}
 =\frac{\log2}{\ell_*}\,\Lambda
\end{equation}
by Eq.~\eqref{ym:scaling-law}. Three further steps are then needed.
\begin{enumerate}
\item \emph{One coupling on both boxes.} $\Delta_\beta(n)$ compares tori at
 the same $\beta$. In Theorem~\ref{ym:thm:existence} the bare coupling
 $\beta_{N,p}$ depends on the volume index $p$ through the condition
 $\beta_{N,p,0}=\bar\beta$ of Eq.~\eqref{ym:running}. The defect bound must
 be stable under the difference $\beta_{N,p}-\beta_{N,p'}$, which
 Eq.~\eqref{ym:reg1} controls only to $o(\log(1/a_N))$. Leading-logarithm
 scaling alone does not bound $a\,n(\beta)$: for example
 $\beta/b_*=\log(1/a)+\log\log(1/a)$ has the asymptotics of
 Eq.~\eqref{ym:scaling-law} while $a\,e^{\beta/b_*}=\log(1/a)\to\infty$.
\item \emph{Passage to infinite volume.} On the spatial torus of side
 $L_k=2^{k+1}n$, let $\Omega_k$ be a normalized eigenvector of $\lambda_1$,
 $\psi_j$ an orthonormal eigenbasis with $\psi_1=\Omega_k$, and $r_k$ the
 ratio $\lambda_2/\lambda_1$ there. For a bounded local observable $A$ of one
 time slice, acting by multiplication, the centred time correlation is
 \begin{equation}\label{ym:finite-correlation}
  C_k(t)=\Big\langle A\Omega_k,\Big(\frac{\mathsf T}{\lambda_1}\Big)^tA\Omega_k
  \Big\rangle-\langle\Omega_k,A\Omega_k\rangle^2
  =\sum_{j\ge2}\Big(\frac{\lambda_j}{\lambda_1}\Big)^t
  \big|\langle\psi_j,A\Omega_k\rangle\big|^2,
 \end{equation}
 where the $j=1$ term of the expansion of $A\Omega_k$ cancels the subtracted
 square. Each ratio is at most $r_k$ and
 $\sum_j|\langle\psi_j,A\Omega_k\rangle|^2=\|A\Omega_k\|^2\le\|A\|^2$, so
 \begin{equation}\label{ym:finite-correlation-bound}
  0\le C_k(t)\le\|A\|^2\,r_k^{\,t}\le\|A\|^2\,2^{-t/n}
 \end{equation}
 by Eq.~\eqref{ym:conditional-gap}. The infinite-volume statement needs a
 hypothesis that is not proved here: for each fixed $t$ the limit
 $C(t)=\lim_{k\to\infty}C_k(t)$ exists and is the centred correlation of $A$
 in the infinite-volume state. Then
 $|C(t)|\le\|A\|^2\,2^{-t/n}$, since the bound
 \eqref{ym:finite-correlation-bound} is uniform in $k$. A gap for the whole
 reconstructed theory needs in addition that such vectors $A\Omega$ be dense,
 control of the vacuum subtraction, the identification of the limit state
 with the one of Theorem~\ref{ym:thm:existence} through the convergence
 \eqref{ym:ordered-limit}, and, with Proposition~\ref{ym:prop:gap-decay},
 the decay for the renormalized insertions with constants uniform in the
 cutoff; then Eq.~\eqref{ym:gap-statement} follows for the reconstructed
 theory.
\item \emph{Uniqueness.} Exponential clustering of $S^{\mathrm{mix}}$ makes
 $X$ a point for that limit. Independence of the subsequence $q$ requires in
 addition that the infinite-volume limit of lattice correlations exists,
 which a uniform gap with exponentially small finite-volume corrections
 would give.
\end{enumerate}
Conjecture~\ref{ym:conj:defect} is the infrared content: by
Eq.~\eqref{ym:lattice-mass} it concerns boxes of side $e^{\beta/b_*}$ in
lattice units, where the effective coupling is of order one.

\subsection{Centre twists}
Suppose $-1\in G$ (for instance $G=SU(2)$, or $SU(N)$ with $N$ even); it is
central. For a set $\Sigma$ of plaquettes define the twisted partition
function by reversing the sign of the plaquette variable on $\Sigma$:
\begin{equation}\label{ym:twisted-Z}
 Z_\beta[\Sigma]=\int\exp\Big(-\beta\sum_{p\notin\Sigma}
 \Big(1-\frac1m\operatorname{Re}\operatorname{tr}U_p\Big)
 -\beta\sum_{p\in\Sigma}
 \Big(1+\frac1m\operatorname{Re}\operatorname{tr}U_p\Big)\Big)
 \prod_{(y,\mu)}h\big(dU(y,\mu)\big).
\end{equation}
Since $1+\frac1m\operatorname{Re}\operatorname{tr}U_p
=\big(1-\frac1m\operatorname{Re}\operatorname{tr}U_p\big)
+\frac2m\operatorname{Re}\operatorname{tr}U_p$,
\begin{equation}\label{ym:twist-ratio}
 \frac{Z_\beta[\Sigma]}{Z_\beta[\emptyset]}
 =\Big\langle\exp\Big(-\frac{2\beta}{m}\sum_{p\in\Sigma}
 \operatorname{Re}\operatorname{tr}U_p\Big)\Big\rangle_{\beta},
\end{equation}
the expectation in the Wilson measure \eqref{ym:wilson-measure}.
\begin{lemma}[Moving a twist]\label{ym:lem:twist-move}
For a link $b$ let $\delta b$ be the set of plaquettes containing $b$. Then
$Z_\beta[\Sigma\,\triangle\,\delta b]=Z_\beta[\Sigma]$, where $\triangle$ is
the symmetric difference.
\end{lemma}
\begin{proof}
Substitute $U(b)\to-U(b)$. Haar measure is invariant under multiplication by
the group element $-1$. By Eq.~\eqref{ym:plaquette} a plaquette variable
contains $U(b)$ or $U(b)^{-1}$ exactly once if $p\in\delta b$ and not at all
otherwise, and $-1$ is central with $(-1)^{-1}=-1$; so $U_p\to-U_p$ for
$p\in\delta b$ and $U_p\to U_p$ otherwise. This exchanges the two summands of
Eq.~\eqref{ym:twisted-Z} exactly for the plaquettes in $\delta b$.
\end{proof}
So $Z_\beta[\Sigma]$ depends only on the class of $\Sigma$ modulo the sets
$\delta b$. On the torus with sides $(n_1,n_2,n_3,n_4)$ let
\begin{equation}\label{ym:thooft-surface}
 \Sigma_{12}=\{\text{plaquettes in the plane }(1,2)\text{ at sites }y
 \text{ with }y^1=y^2=0\}.
\end{equation}
\begin{lemma}\label{ym:lem:thooft-class}
Every 3-cube contains an even number of plaquettes of $\Sigma_{12}$, and
$\Sigma_{12}$ is not a symmetric difference of sets $\delta b$.
\end{lemma}
\begin{proof}
A cube in the directions $(1,2,3)$ at $y$ contains two plaquettes in the
plane $(1,2)$, at $y$ and $y+e_3$; both or neither have $y^1=y^2=0$. The same
holds for $(1,2,4)$. Cubes in the directions $(1,3,4)$ and $(2,3,4)$ contain
no plaquette in the plane $(1,2)$.

Fix $(y^3,y^4)$ and let $\Pi$ be the set of all plaquettes in the plane
$(1,2)$ at that $(y^3,y^4)$. Then $|\Sigma_{12}\cap\Pi|=1$. A link in
direction $1$ or $2$ at that $(y^3,y^4)$ lies in exactly two plaquettes of
$\Pi$; any other link lies in none. So $|\delta b\cap\Pi|$ is even for every
$b$, and so is $|S\cap\Pi|$ for every symmetric difference $S$ of sets
$\delta b$.
\end{proof}
The first property says that the twist creates no monopole: the product of
the signs over the faces of each cube is $+1$. The twist free energy is
\begin{equation}\label{ym:twist-free-energy}
 F_\beta(n_1,n_2;n_3,n_4)=-\log\frac{Z_\beta[\Sigma_{12}]}{Z_\beta[\emptyset]}.
\end{equation}
\paragraph{Example: $G=\{\pm1\}$ in two dimensions.}
Take $G=\{\pm1\}\subset U(1)$, $m=1$, on the two-dimensional torus with sides
$(L,L)$: $P=L^2$ plaquettes, $2P$ links, each link in exactly two plaquettes.
Write $\sigma_p=U_p=\pm1$ and $s_p=-1$ for $p\in\Sigma$, $s_p=+1$ otherwise.
For $\sigma=\pm1$, $e^{\beta\sigma}=\cosh\beta+\sigma\sinh\beta$, so
Eq.~\eqref{ym:twisted-Z} with normalized counting measure on the links is
\begin{align}
 Z_\beta[\Sigma]
 &=e^{-\beta P}\,2^{-2P}\sum_{\sigma}\prod_p
   \big(\cosh\beta+s_p\sigma_p\sinh\beta\big)\notag\\
 &=e^{-\beta P}\sum_{S}\cosh^{P-|S|}\beta\,\sinh^{|S|}\beta\,
   \Big(\prod_{p\in S}s_p\Big)\,
   2^{-2P}\sum_{\sigma}\prod_{p\in S}\sigma_p ,
 \label{ym:z2-expansion}
\end{align}
the sum over all sets $S$ of plaquettes. Let $n_b(S)\in\{0,1,2\}$ be the
number of plaquettes of $S$ containing the link $b$. Then
$\prod_{p\in S}\sigma_p=\prod_b\sigma_b^{\,n_b(S)}$, and the average over
$\sigma_b=\pm1$ vanishes unless every $n_b(S)$ is even, in which case it is
$1$. $n_b(S)$ even for every $b$ means that the two plaquettes at each link
are both in $S$ or both outside; adjacency of plaquettes across links is
connected, so $S=\emptyset$ or $S$ is the set of all plaquettes:
\begin{equation}\label{ym:z2-closed}
 Z_\beta[\Sigma]=e^{-\beta P}\Big(\cosh^P\beta+(-1)^{|\Sigma|}\sinh^P\beta\Big),
 \qquad
 \frac{Z_\beta[\Sigma]}{Z_\beta[\emptyset]}
 =\frac{1-\tanh^{L^2}\beta}{1+\tanh^{L^2}\beta}\quad(|\Sigma|\text{ odd}).
\end{equation}
The twist free energy
$\log\frac{1+\tanh^{L^2}\beta}{1-\tanh^{L^2}\beta}$ tends to $0$ as
$L\to\infty$ for every $\beta$, like $2\tanh^{L^2}\beta$: exponentially in the
area. The dependence on $\Sigma$ only through the parity of $|\Sigma|$ is
Lemma~\ref{ym:lem:twist-move}: here $\delta b$ consists of two plaquettes.

In the two-dimensional $O(n)$ models, $n\ge3$, the analogous ratio for a
rotation by $\pi$ across an annulus tends to $1$ as the annulus grows, at
every temperature, and this drives the exponential decay of correlations
\cite{OpenAIOn2026}.
\begin{conjecture}[Vanishing twist free energy]\label{ym:conj:twist}
For $G=SU(2)$ and every $\beta\ge0$,
$F_\beta(n,n;n,n)\to0$ as $n\to\infty$.
\end{conjecture}
The relation between Conjecture~\ref{ym:conj:twist}, with a rate, and
Conjecture~\ref{ym:conj:defect} is open. Two
differences from the spin model are already visible in
Eq.~\eqref{ym:twist-ratio}. The signs of $\operatorname{tr}U_p$ are not
those of a ferromagnetic system: a configuration can have an odd number of
negative plaquettes on the faces of a cube. And the twist is carried by a
surface, so the planar crossing arguments have no direct counterpart.

\subsection{Positivity of the transfer operator}\label{ym:sc-positivity}
Proposition~\ref{ym:prop:defect} uses that the operator $\mathsf T$ of
Eq.~\eqref{ym:transfer-kernel} is nonnegative on every spatial torus. This
holds for every compact $G\subset U(m)$ and every $\beta\ge0$; the proof below
uses only inner products and products of kernels.

A real kernel $k\colon X\times X\to\mathbb R$ is \emph{positive semidefinite}
if it is symmetric, $k(x,x')=k(x',x)$, and $\sum_{i,j}c_ic_j\,k(x_i,x_j)\ge0$
for all finite families of points $x_i\in X$ and reals $c_i$.
\begin{lemma}[Products and exponentials]\label{ym:sc-lem:schur}
\begin{enumerate}
\item If $\phi\colon X\to\mathbb R^D$, then $k(x,x')=\phi(x)\cdot\phi(x')$ is
 positive semidefinite.
\item If $k_1,k_2$ are positive semidefinite on $X$, so is the pointwise
 product $k_1k_2$; if $k_1$ is positive semidefinite on $X$ and $k_2$ on $X'$,
 then $k_1(x,x')\,k_2(y,y')$ is positive semidefinite on $X\times X'$.
\item If $k$ is positive semidefinite and $\beta\ge0$, then $e^{\beta k}$ is
 positive semidefinite.
\end{enumerate}
\end{lemma}
\begin{proof}
Each kernel constructed is symmetric because its factors are. (1)
$\sum_{i,j}c_ic_j\,\phi(x_i)\cdot\phi(x_j)=\big|\sum_ic_i\phi(x_i)\big|^2$.
(2) For finitely many points the matrices $A_{ij}=k_1(x_i,x_j)$ and
$B_{ij}=k_2(x_i,x_j)$ are positive semidefinite. Write
$B=\sum_r b_rb_r^{\mathsf T}$ with vectors $b_r$. Then
$\sum_{i,j}c_ic_jA_{ij}B_{ij}=\sum_r\sum_{i,j}(c_ib_{r,i})(c_jb_{r,j})A_{ij}
\ge0$ (Schur's product theorem). The second statement is the first applied on
$X\times X'$ to the kernels $k_1(x,x')$ and $k_2(y,y')$, each positive
semidefinite there. (3) By (2) every power $k^r$ is positive semidefinite, and
$e^{\beta k}=\sum_{r\ge0}\frac{\beta^r}{r!}k^r$ is a pointwise limit of
nonnegative combinations of them.
\end{proof}
On $m\times m$ complex matrices, $\langle A,B\rangle=\operatorname{Re}
\operatorname{tr}(AB^\dagger)$ is the Euclidean inner product of
$\mathbb R^{2m^2}$. For unitary $U'$, $U'^{-1}=U'^\dagger$, so the single-link
factor of the temporal weight,
\begin{equation}\label{ym:sc-link-kernel}
 k(U,U')=\exp\Big(-\beta\Big(1-\frac1m\operatorname{Re}\operatorname{tr}\,
 UU'^{-1}\Big)\Big)=e^{-\beta}\exp\Big(\frac\beta m\langle U,U'\rangle\Big),
\end{equation}
is positive semidefinite on $G$ by Lemma~\ref{ym:sc-lem:schur}. For
$G=SU(2)$, $U=u^0+iu^k\sigma_k$ with $u\in S^3\subset\mathbb R^4$, and
$\frac12\operatorname{Re}\operatorname{tr}UU'^{-1}=u\cdot u'$.
\begin{proposition}\label{ym:sc-prop:positivity}
Let $G\subset U(m)$ be compact and $\beta\ge0$. On every spatial torus the
operator $\mathsf T$ of Eq.~\eqref{ym:transfer-kernel} on $L^2$ of the spatial
configurations, with product Haar measure, is self-adjoint, nonnegative and
trace class, and $\operatorname{Tr}\mathsf T>0$.
\end{proposition}
\begin{proof}
Write $\mathcal U_{\mathrm s}$ for the spatial configurations of one slice.
For a gauge transformation $\gamma\colon\mathbf x\mapsto\gamma(\mathbf x)\in G$
on the slice, $U^\gamma$ is given by Eq.~\eqref{ym:link-gauge}, and
$(R_\gamma f)(U)=f(U^\gamma)$ is unitary on $L^2(\mathcal U_{\mathrm s})$,
because Haar measure is invariant under left and right multiplication.

\emph{The gauge average.} In Eq.~\eqref{ym:transfer-kernel}, put
$\gamma=V$. Then
$V(\mathbf x)U'(\mathbf x,j)V(\mathbf x+e_j)^{-1}=U'^{V}(\mathbf x,j)$ and
\begin{equation}
 \operatorname{Re}\operatorname{tr}\,U(\mathbf x,j)V(\mathbf x+e_j)
 U'(\mathbf x,j)^{-1}V(\mathbf x)^{-1}
 =\operatorname{Re}\operatorname{tr}\,U(\mathbf x,j)\,
 U'^{V}(\mathbf x,j)^{-1}.
\end{equation}
Hence, with $\mathsf K_0(U,U'')=\prod_{\mathbf x,j}k\big(U(\mathbf x,j),
U''(\mathbf x,j)\big)$,
\begin{equation}\label{ym:sc-gauge-average}
 \mathsf K(U,U')=\int\prod_{\mathbf x}h\big(dV(\mathbf x)\big)\,
 \mathsf K_0\big(U,U'^{V}\big).
\end{equation}
Since $U^\gamma(U''^\gamma)^{-1}=\gamma(\mathbf x)\,UU''^{-1}\gamma(\mathbf
x)^{-1}$ link by link and the trace is invariant under conjugation,
$\mathsf K_0(U^\gamma,U''^\gamma)=\mathsf K_0(U,U'')$: the operator
$\mathsf K_0$ commutes with every $R_\gamma$. Substituting $U''=U'^V$ in
$\int\mathsf K(U,U')f(U')\,dU'$, which preserves the measure, gives
$\mathsf K=\mathsf K_0\mathsf P$ with
$\mathsf P=\int\prod_{\mathbf x}h(d\gamma(\mathbf x))\,R_\gamma$, the orthogonal
projection onto gauge-invariant functions. As $\mathsf P$ commutes with
$\mathsf K_0$, $\mathsf K=\mathsf P\mathsf K_0\mathsf P$.

\emph{Positivity.} By Lemma~\ref{ym:sc-lem:schur}(2) the kernel $\mathsf K_0$ is
positive semidefinite on $\mathcal U_{\mathrm s}$. It is continuous, and
$\mathcal U_{\mathrm s}$ is compact with a finite measure. Approximating
$\int\int f(U)\mathsf K_0(U,U')f(U')\,dU\,dU'$ by Riemann sums over finite
partitions shows $\langle f,\mathsf K_0f\rangle\ge0$ for continuous $f$, and by
density for all $f\in L^2$. Hence
$\langle f,\mathsf Kf\rangle=\langle\mathsf Pf,\mathsf K_0\mathsf Pf\rangle\ge0$.
With $\mathsf M$ the multiplication by $e^{-\frac\beta2E_{\mathrm s}}$, a
bounded positive self-adjoint operator, $\mathsf T=\mathsf M\mathsf K\mathsf M$
and $\langle f,\mathsf Tf\rangle=\langle\mathsf Mf,\mathsf K\mathsf Mf\rangle
\ge0$.

\emph{Trace class.} By Mercer's theorem a continuous positive semidefinite
kernel on a compact space with a finite Borel measure defines a trace-class
operator whose trace is the integral of the diagonal \cite{Mercer1909}. For
$\mathsf K_0$ the diagonal is $\prod k(U,U)=1$, so
$\operatorname{Tr}\mathsf K_0=1$. Products of a trace-class operator with
bounded operators are trace class, so $\mathsf T=\mathsf M\mathsf P\mathsf
K_0\mathsf P\mathsf M$ is trace class. Finally
$\operatorname{Tr}\mathsf T=Z(n_1,n_2,n_3,1)>0$ by
Eq.~\eqref{ym:trace-formula}.
\end{proof}
The Wilson energy is symmetric under permutations of the four axes, so the
same holds with any axis as time, and the hypothesis of
Proposition~\ref{ym:prop:defect} holds for every rectangular torus. This is
the positivity of the Wilson transfer matrix \cite{Luscher1977,
OsterwalderSeiler1978}; the argument above replaces the character expansion
by Lemma~\ref{ym:sc-lem:schur}. For a finite group such as
$G=\{\pm1\}$, $L^2(\mathcal U_{\mathrm s})$ is finite-dimensional, $\mathsf T$
is a positive semidefinite matrix, and Mercer's theorem is not needed.

\subsection{Measurements of the defect}\label{ym:num-defect}
The partition function of Eq.~\eqref{ym:wilson-partition} satisfies
$Z(\beta=0)=1$ and
$\frac{\partial\log Z}{\partial\beta}=-\langle E\rangle_\beta$. On a torus with
$N_P$ plaquettes write
$\langle a_0\rangle=\frac{1}{N_P}\sum_p\frac1m\langle\operatorname{Re}
\operatorname{tr}U_p\rangle$. Since $N_P(2n,2n,2n,2n)=16\,N_P(n,n,n,n)$, the
constant parts of the energy cancel in Eq.~\eqref{ym:defect}, and
\begin{equation}\label{ym:num-integration}
 \Delta_\beta(n)=N_P(2n,2n,2n,2n)\int_0^\beta\Big(\langle a_0\rangle_{n^4}
 -\langle a_0\rangle_{(2n)^4}\Big)\,d\beta'.
\end{equation}
The defect is the large volume times the difference of plaquette averages
between the torus and its double. Each $\log\mathsf P_k$ of
Eq.~\eqref{ym:four-purities} has the same form between consecutive boxes.
Simulations use heat-bath and over-relaxation updates for $G=SU(2)$, with
Simpson integration on a grid of step $0.1$ in $\beta$ and jackknife errors.
The code is \texttt{Rust/lattice-gauge}, with CUDA C++ updates.

\paragraph{Check in two dimensions.} For $SU(2)$ on a two-dimensional torus with
$V$ plaquettes, $Z=\sum_{j}\big(\frac{2I_{2j+1}(\beta)}{\beta}\big)^V$, normalized
as in Eq.~\eqref{ym:num-integration}. The square defect
$4\log Z(n,n)-\log Z(2n,2n)$ at $n=2$:
\begin{center}
\begin{tabular}{cccccc}
$\beta$ & 1.2 & 2.0 & 2.4 & 3.2 & 4.0\\\hline
Monte Carlo & 0.0204(49) & 0.1398(62) & 0.2342(66) & 0.4746(72) & 0.7375(77)\\
exact & 0.0258 & 0.1396 & 0.2343 & 0.4718 & 0.7335
\end{tabular}
\end{center}

\paragraph{$SU(2)$, $n=2$.} Two replicas of $4\times10^5$ updates per point
and box:
\begin{center}
\begin{tabular}{ccccccc}
$\beta$ & $\Delta_\beta(2)$ & $\log\mathsf P_1$ & $\log\mathsf P_2$ &
 $\log\mathsf P_3$ & $\log\mathsf P_4$\\\hline
0.4 & 0.056(18) & $-0.002(3)$ & $-0.010(4)$ & $-0.006(5)$ & 0.012(8)\\
0.6 & 0.187(23) & $-0.010(3)$ & $-0.022(5)$ & $-0.015(7)$ & 0.010(9)\\
0.8 & 0.591(27) & $-0.037(4)$ & $-0.053(6)$ & $-0.047(8)$ & 0.005(11)\\
1.0 & 1.456(31) & $-0.093(5)$ & $-0.120(6)$ & $-0.113(9)$ & $-0.006(12)$\\
1.6 & 13.14(4) & $-0.869(6)$ & $-1.019(9)$ & $-0.934(12)$ & $-0.248(16)$\\
2.0 & 43.61(6) & $-2.430(8)$ & $-2.959(12)$ & $-3.428(16)$ & $-5.474(22)$\\
2.4 & 69.54(6) & $-3.309(10)$ & $-3.861(14)$ & $-6.067(19)$ & $-15.49(3)$
\end{tabular}
\end{center}
The inequalities $\mathsf P_k\le1$ and $\mathsf P_4\ge e^{-\Delta}$ of
Proposition~\ref{ym:prop:defect} hold within errors at every $\beta$. The rapid
growth between $\beta=1$ and $2$ lies near the coupling of the $N_t=2$
thermal transition; a finite symmetric torus has no phase transition, so it is
a finite-size crossover.

\paragraph{$SU(2)$, $n=4$.} The $8^4$ scans start at $\beta_0=1$, so the
measured quantity is the increment $\Delta_\beta(4)-\Delta_{1}(4)$:
\begin{center}
\begin{tabular}{ccccccc}
$\beta$ & 1.6 & 1.8 & 2.0 & 2.2 & 2.4 & 2.6\\\hline
$\Delta_\beta(4)-\Delta_1(4)$ & $-0.012(41)$ & $-0.003(51)$ & 0.336(63) &
 9.32(8) & 40.89(9) & 53.43(10)
\end{tabular}
\end{center}

\paragraph{Strong coupling.} Expand each plaquette weight in characters,
$e^{\frac\beta2\operatorname{tr}U}=c_0(\beta)\big(1+\sum_{j\ge\frac12}d_j\,
u_j(\beta)\chi_j(U)\big)$ with $d_j=2j+1$ and $u_{1/2}=u=\frac{I_2(\beta)}
{I_1(\beta)}$. Then $\log Z$ is a sum over connected clusters of plaquettes.
An embedded orientable closed surface of $F$ fundamental plaquettes, $E$ links
and $V$ sites carries $2^{F-E+V}u^F$: a factor $d_{1/2}=2$ per plaquette from
the expansion, $\frac12$ per link from the integration
$\int\chi_{1/2}(UW)\chi_{1/2}(U^{-1}W')\,h(dU)=\frac12\chi_{1/2}(WW')$, and $2$
per closed index loop at a site. Clusters that lift with the same weight and
density to the doubled torus cancel in $\Delta$. The leading survivors are
surfaces of Euler characteristic $0$ that wind around the torus: the $6n^2$
flat tori in coordinate planes, of area $n^2$, and the $12n^3$ square tubes
around one axis, of area $4n$. With
$S(n)=6n^2u^{n^2}+12n^3u^{4n}$ these give $16\,S(n)-S(2n)$, so
\begin{equation}\label{ym:num-strong-coupling}
 \Delta_\beta(n)\approx16\big(6n^2u^{n^2}+12n^3u^{4n}\big).
\end{equation}
The approximation is a truncation, not a controlled expansion: at $n=2$ a
flat sheet in the representation $j=1$ has weight $(I_3/I_1)^4=O(\beta^8)$, the
order of the tube, and the general remainder is not bounded here. Against the
measurements:
\begin{center}
\begin{tabular}{cccc}
$n$, $\beta$ & $u$ & Eq.~\eqref{ym:num-strong-coupling} & measured\\\hline
2, 0.4 & 0.099 & 0.037 & 0.056(18)\\
2, 0.6 & 0.148 & 0.184 & 0.187(23)\\
2, 0.8 & 0.195 & 0.557 & 0.591(27)\\
2, 1.0 & 0.240 & 1.295 & 1.456(31)
\end{tabular}
\end{center}
Equation~\eqref{ym:num-strong-coupling} equals $\frac1{7480}$ at
$\beta=0.0972$ for $n=2$ and $\beta=1.355$ for $n=4$.

\paragraph{The abelian counter-model.} For $G=U(1)$ the same construction
applies, and Proposition~\ref{ym:sc-prop:positivity} holds. At large $\beta$
the four-dimensional theory has a massless photon. The measured defect,
with the $8^4$ values as increments from $\beta_0=0.6$:
\begin{center}
\begin{tabular}{ccccc}
$\beta$ & 0.9 & 1.0 & 1.2 & 1.6\\\hline
$\Delta_\beta(2)$ & 50.23(2) & 71.30(3) & 79.56(3) & 88.48(3)\\
$\Delta_\beta(4)-\Delta_{0.6}(4)$ & 0.17(5) & 34.3(1) & 70.8(1) & 80.3(1)
\end{tabular}
\end{center}
In the Coulomb phase the defect does not decrease when the box doubles, and
the integrand of Eq.~\eqref{ym:num-integration} is the same at $n=2$ and
$n=4$ within $3\%$ ($17.84(4)$ and $18.3(2)$ at $\beta=1.6$). For $SU(2)$ at
$\beta=1.6$--$2.0$, $\Delta_\beta(2)$ grows from $13$ to $44$ while
$\Delta_\beta(4)-\Delta_1(4)$ is zero within errors up to $\beta=1.8$. A gapless theory admits no seed $\Delta_\beta(n)\le\frac1{7480}$,
once the finite-volume bound of Proposition~\ref{ym:prop:defect} is carried to
the infinite-volume state, so this is where an argument must use the
non-abelian group.

\subsection{The \texorpdfstring{$0^{++}$}{0++} glueball against the size of the box}\label{ym:num-glueball}
On $L^3\times T$ with time along axis 4, the zero-momentum operators
$O_i(t)$ are the averages of $\frac12\operatorname{tr}U_p$ over spatial
plaquettes of slice $t$, after $0,4,8,16$ APE smearing steps of the spatial
links. With $C_{ij}(t)=\langle O_i(t_0)O_j(t_0+t)\rangle-\langle O_i\rangle
\langle O_j\rangle$, let $v$ be the eigenvector of $C(1)v=\lambda C(0)v$ with
the largest $\lambda$, and $c(t)=v^{\mathsf T}C(t)v$. A one-state fit
$c(t)=A\big(e^{-mt}+e^{-m(T-t)}\big)$ is made over each window, with the
jackknife covariance. The window $[2,5]$ is used if its $\chi^2$ per degree of
freedom is below $2$, otherwise $[3,6]$ under the same condition. At $\beta=2.5$
(and the reference values of \cite{Teper1998}, Table~10):
\begin{center}
\begin{tabular}{cccc}
$L^3\times T$ & $am$, $[2,5]$ ($\chi^2$/dof) & $am$, $[3,6]$ ($\chi^2$/dof) &
 \cite{Teper1998}\\\hline
$8^3\times16$ & 0.714(7) (31) & 0.630(12) (11) & \\
$10^3\times16$ & 0.624(7) (31) & 0.565(9) (7.0) & \\
$12^3\times16$ & 0.677(7) (11) & 0.636(12) (3.8) & \\
$14^3\times16$ & 0.726(8) (3.6) & 0.687(18) (1.2) & \\
$16^3\times16$ & 0.725(9) (4.4) & 0.691(17) (1.8) & 0.660(22)\\
$20^3\times20$ & 0.730(11) (0.43) & 0.711(22) (0.02) & 0.723(23)
\end{tabular}
\end{center}
At $\beta=2.4$ on $16^3\times16$ the window $[2,5]$ gives $0.927(30)$ against
$0.936(23)$. These are comparisons at matched volumes; they do not validate
the one-state fit model. For $L\le12$ both windows fail, and the periodic
effective masses keep decreasing with $t$.

A flux tube winding once around a spatial direction of length $L$ (a
torelon) is created by the Polyakov loop
$\frac12\operatorname{tr}\prod_{k=0}^{L-1}U(\mathbf x+ke_j,j)$. Its energy
$E_T(L)$ is measured from the zero-momentum correlator of that operator,
averaged over the three spatial directions, with the same smearing, fit and
rule. To leading order in $\frac1L$, $E_T(L)=\sigma L-\frac{\pi}{3L}$ with
$\sigma$ the string tension. A single winding loop carries centre charge and
does not mix with the plaquettes. A pair winding in opposite directions is
centre neutral and has zero momentum, so it can mix with the scalar channel,
with energy near $2E_T(L)$ if the two strings do not interact:
\begin{center}
\begin{tabular}{cccccc}
$L$ & 8 & 10 & 12 & 14 & 16\\\hline
$aE_T(L)$ & 0.1925(9) & 0.2368(13) & 0.3123(22) & 0.4023(27) & 0.4851(42)\\
$2aE_T(L)$ & 0.385 & 0.474 & 0.625 & 0.805 & 0.970\\
$a\sqrt\sigma$ & 0.201 & 0.185 & 0.183 & 0.185 & 0.186
\end{tabular}
\end{center}
At $L=20$, $aE_T=0.660(17)$. For $L\ge10$, $a\sqrt\sigma=0.183$--$0.189$;
\cite{Teper1998}, Table~10, gives $0.1834(26)$ on $16^4$ and $0.1881(28)$ on
$20^4$. The value $2aE_T$ lies below the glueball for $L\le12$ and above it
from $L=14$, and at $L=10$ the plaquette correlator over $t\in[5,8]$ (a window
chosen after the fact) gives $0.484(17)$, equal to $2aE_T$ within errors. That
the failing fits at $L\le12$ come from this mixing is a hypothesis: the pair
operator itself was not measured.

No defect is measured here. The chain gives a necessary condition for a seed.
In the telescope of Eq.~\eqref{ym:defect-telescope} at side $n$, take time
as the first doubled direction, with weight $8$ in
$\Delta_\beta(n)=-\log(\mathsf P_1^{\,8}\mathsf P_2^{\,4}\mathsf P_3^{\,2}
\mathsf P_4)$, so $-8\log\mathsf P_1\le\Delta_\beta(n)$. By
Eq.~\eqref{ym:purity-bounds},
$e^{-n\mathsf m(n)}\le s_n\le\frac1{\mathsf P_1}-1$, where $\mathsf m(n)$ is
the transfer gap on the spatial torus of side $n$. Hence
$\Delta_\beta(n)\le\frac1{7480}$ requires
\begin{equation}\label{ym:num-necessary}
 n\,\mathsf m(n)\ \ge\ -\log\big(e^{1/59840}-1\big)=10.9994:
\end{equation}
a seed box is at least about eleven of its own correlation lengths wide. A
certified upper bound on $\mathsf m(n)$ below $\frac{10.9994}{n}$ would exclude
that $n$; effective masses from finite fits are not such a certificate.

\subsection{Topology of centre vortices and monopoles}\label{ym:num-topology}
\paragraph{Vortex sheets.} After direct maximal centre gauge
\cite{DelDebbioEtAl1997} and projection $U\mapsto\operatorname{sign}
\frac12\operatorname{tr}U$, the plaquettes with value $-1$ are dual to closed
surfaces on the dual lattice (the dual P-vortex sheet), since the projected
plaquette field is a coboundary. For a finite 2-complex $X$, $H_2(X;\mathbb Z)$
is free, and by universal coefficients
\begin{equation}\label{ym:num-betti}
 b_2(X;\mathbb F_p)=\operatorname{rank}H_2(X;\mathbb Z)+t_p(H_1(X;\mathbb Z)),
\end{equation}
where $t_p$ counts the cyclic factors of $p$-power order. The difference
$b_2(\mathbb F_2)-b_2(\mathbb F_p)$, $p=1\,000\,003$, is $t_2-t_p$; a value
$1$ proves at least one factor of even order in $H_1$ of the whole sheet. The
number of configurations, of $20$ (of $10$ for the entry marked $*$), with
difference $1$:
\begin{center}
\begin{tabular}{cccc}
$\beta$ & $8^4$ & $12^4$ & $16^4$\\\hline
2.2 & 20 & & \\
2.3 & 20 & & \\
2.4 & $10^*$ & & \\
2.5 & 7 & 20 & \\
2.6 & 4 & 14 & 19\\
2.7 & & & 15\\
2.8 & 1 & & 1\\
2.9 & & & 3
\end{tabular}
\end{center}
Random sheets at the same density on $8^4$ give difference $0$ in four of four
samples. On $16^4$ the drop lies between $\beta=2.7$ and $2.8$, around the
$N_t=16$ thermal coupling $2.7395(100)$ \cite{Teper1998}, Table~14. This is an
association with confinement, for one gauge-fixing prescription. It is not an
infinite-volume order parameter.

\paragraph{Monopole currents.} After maximal abelian gauge the
DeGrand--Toussaint currents \cite{DeGrandToussaint1980} are integer 1-cycles.
The fraction of the current in the largest cluster on $12^4$ is
$0.55,0.35,0.18,0.15$ at $\beta=2.4$--$2.7$, below one half already at
$\beta=2.5$, where the vortex count is still $20$ of $20$. For $U(1)$ on $6^4$
it is $0.99$ at $\beta=0.8$ (confined) and $0.22$ at $\beta=1.2$ (Coulomb).
